\section{Secant objects on the split member}\label{sec:secantobjects}

\subsection{Sums of line bundles and the positivity obstruction}
\label{ssec:whatobject}

\Cref{thm:semiregclosure} asks for one perfect complex at one base point.
The cheapest candidate is a sum of line bundles on a split member, because
there the Weil classes are already polynomials in $\beta$, $\widehat\beta$ and
$\ell$ by \Cref{thm:splitclosed}, so the Chern character one wants is
manifestly in reach. This subsection settles that candidate. It does not
exist, and the obstruction is not semiregularity: it is the Chern character
condition, and it fails because the norm form of $K$ is positive
definite.

The norm rules out three different constructions in this paper, by three
different mechanisms, and they should be kept apart.
\Cref{thm:character} uses the norm \emph{character} of $\Kx$ and an
equivariance hypothesis. \Cref{prop:divfails} uses the fact that the part $R$
of $H^{2}$ on which $\Kx$ acts through $\tau^{2}$ and $\bbar\tau^{\,2}$ has no
Hodge classes at a general point, which is a statement about the Hodge group.
\Cref{thm:positivity} below uses neither: it uses that $\Nm_{K/\QQ}$ is a
positive definite quadratic form on $K$ and that the multiplicities of a
filtration are positive integers. The three are independent, and only the
third is available at a special point, which is exactly where the object of
\Cref{prop:ktheory} has to live.

\begin{lemma}[The semiregularity map of a split object]\label{lem:splitsemireg}
Let $Y$ be a smooth projective complex variety and let
\[
  E=\bigoplus_{i=1}^{s}L_{i}[p_{i}]^{\oplus m_{i}},
  \qquad \lambda_{i}=c_{1}(L_{i}),
\]
with the $L_{i}$ line bundles, the $\lambda_{i}$ pairwise distinct and
$p_{i}\in\ZZ$. Then for $\zeta=(\zeta_{ij})\in\Ext^{2}(E,E)$,
\begin{equation}\label{eq:splitsigma}
  \sigma_{E}(\zeta)_{q}
  =\frac{1}{q!}\sum_{i=1}^{s}(-1)^{p_{i}}\,
    \mathrm{tr}(\zeta_{ii})\,\lambda_{i}^{\,q},
\end{equation}
so $\sigma_{E}$ annihilates $\Ext^{2}(L_{i}[p_{i}],L_{j}[p_{j}])$ for $i\ne j$
and the traceless part of each diagonal block. Consequently $\sigma_{E}$ is
injective if and only if all three of the following hold:
\begin{enumerate}[label=\textup{(\roman*)},nosep]
\item $m_{i}=1$ for every $i$;
\item $H^{\,2+p_{i}-p_{j}}\bigl(Y,L_{j}L_{i}^{-1}\bigr)=0$ for all $i\ne j$;
\item the map
  \[
    \Phi\colon H^{2}(\cO_{Y})^{\oplus s}\longrightarrow H^{*}(Y,\CC),
    \qquad
    (\zeta_{i})\longmapsto\sum_{i}e^{\lambda_{i}}\zeta_{i},
  \]
  is injective.
\end{enumerate}
\end{lemma}

\begin{proof}
The Atiyah class of a direct sum is the direct sum of the Atiyah classes, and
that of a line bundle is its first Chern class, so $\mathrm{At}(E)$ is
diagonal with entries $\lambda_{i}$ and $\mathrm{At}(E)^{q}$ is diagonal with
entries $\lambda_{i}^{\,q}$. The matrix product $\zeta\circ\mathrm{At}(E)^{q}$
therefore has $(i,j)$ entry $\zeta_{ij}\lambda_{i}^{\,q}$, and the trace keeps
only the diagonal; a shift by $p$ multiplies the trace by $(-1)^{p}$. That is
\eqref{eq:splitsigma}, and it exhibits the two spaces the map kills.

For the equivalence, the off-diagonal part of $\Ext^{2}(E,E)$ is
$\bigoplus_{i\ne j}\Ext^{2}(L_{i}[p_{i}],L_{j}[p_{j}])
=\bigoplus_{i\ne j}H^{2+p_{i}-p_{j}}(L_{j}L_{i}^{-1})$, which must vanish,
giving (ii); the diagonal block at $i$ is $H^{2}(\cO_{Y})\otimes
\mathrm{End}(\CC^{m_{i}})$ and only its trace survives, so $m_{i}=1$, giving
(i). What is left is $\bigoplus_{i}H^{2}(\cO_{Y})$, and by
\eqref{eq:splitsigma} an element $(\zeta_{i})$ of it dies exactly when
$\sum_{i}\lambda_{i}^{\,q}\zeta_{i}=0$ for every $q\ge0$, which is (iii) after
absorbing the signs $(-1)^{p_{i}}$ into the $\zeta_{i}$ and the factorials
into the grading.
\end{proof}

Condition (ii) is cheap, because the shifts $p_{i}$ are at our disposal: only
their parities are fixed by the Chern character, and a line bundle whose first
Chern class is nondegenerate of index $k$ has cohomology in degree $k$ alone,
so one may choose the $p_{i}$ to avoid every index that occurs. Condition
(iii) is a rank condition, and condition (i) is free. What is not free is the
Chern character itself, and that is the subject of the rest of this
subsection.

\begin{notation}[Coordinates on the divisor classes]\label{not:splitcoords}
Let $A_{0}=X\times\widehat{X}$ be a split member, $n\ge2$, and keep
$\gamma=d\beta-\widehat{\beta}$ from \Cref{not:threeclasses}. Write
\[
  \theta^{\pm}=\gamma\mp\sqrt{-d}\,\ell\in\textstyle\bigwedge^{2}V_{\pm},
  \qquad
  R_{0}=\QQ\gamma\oplus\QQ\ell\subset R .
\]
Every $\theta\in R_{0}$ is $\theta=u\theta^{+}+\bbar u\,\theta^{-}$ for a
unique $u\in K$, and in coordinates $u=g+h\sqrt{-d}$ corresponds to
$\theta=2g\gamma+2dh\,\ell$. By \Cref{prop:splitgeom}(iv) the group
$\NS(A_{0})_{\QQ}$ at a split member with $X$ general is
$\QQ\eta\oplus R_{0}$, so every divisor class $\lambda$ is $e\eta+\theta$ with
$e\in\QQ$ and $\theta\in R_{0}$, and we write $u(\lambda)\in K$ for the
coordinate of its $R_{0}$-part.
\end{notation}

\begin{lemma}[A basis in degree four]\label{lem:sixmonomials}
For $n\ge2$ the six products
$\beta^{2},\widehat\beta^{2},\ell^{2},\beta\widehat\beta,\beta\ell,
\widehat\beta\ell$
are linearly independent in $H^{4}(A_{0},\QQ)$. Equivalently the six classes
$\eta^{2}$, $\eta\theta^{+}$, $\eta\theta^{-}$, $(\theta^{+})^{2}$,
$(\theta^{-})^{2}$, $\theta^{+}\theta^{-}$ are a basis of the degree four part
of the subring generated by the divisor classes.
\end{lemma}

\begin{proof}
Grade $H^{*}(A_{0},\QQ)=\bigwedge^{*}U\otimes\bigwedge^{*}U^{\vee}$ by
bidegree. Then $\beta$ has bidegree $(2,0)$, $\widehat\beta$ has $(0,2)$ and
$\ell$ has $(1,1)$, so the six products have bidegrees $(4,0)$, $(0,4)$,
$(2,2)$, $(2,2)$, $(3,1)$, $(1,3)$. All are distinct except for the pair
$\beta\widehat\beta$ and $\ell^{2}$, and those two are independent in
$\bigwedge^{2}U\otimes\bigwedge^{2}U^{\vee}$ because
$\beta\widehat\beta=\beta\otimes\widehat\beta$ is decomposable there while
$\ell^{2}=-2\sum_{i<j}(x_{i}x_{j})\otimes(\xi_{i}\xi_{j})$ is diagonal in
the induced bases, of tensor rank $\binom{2n}{2}>1$. The second
description follows because the change of basis between the two sets of six
classes is invertible over $K$: $\eta=d\beta+\widehat\beta$ and
$\gamma=d\beta-\widehat\beta$ span the same plane as
$\beta$ and $\widehat\beta$ since $\det\begin{psmallmatrix}d&1\\
d&-1\end{psmallmatrix}=-2d\ne0$, and $\theta^{\pm}$ span the same plane
as $\gamma$ and $\ell$ over $\CC$.
\end{proof}

The next lemma says the same thing at every abelian variety of Weil type, not
only at a split member, and it says it without coordinates. It is the fact on
which the positivity argument turns.

\begin{lemma}[$\theta^{+}\theta^{-}$ is never proportional to $\eta^{2}$]
\label{lem:tensorrank}
Let $(A,\iota,\eta)$ be of $(K,-1,n)$-Weil type with $n\ge2$, and let
$0\ne\theta^{\pm}\in\bigwedge^{2}V_{\pm}$. Then $\theta^{+}\theta^{-}$ is not a
scalar multiple of $\eta^{2}$ in $H^{4}(A,\CC)$.
\end{lemma}

\begin{proof}
Both classes lie in the summand
$\bigwedge^{2}V_{+}\otimes\bigwedge^{2}V_{-}$ of $H^{4}(A,\CC)$, the isotypic
component for the character $\Nm^{2}$, so the question is one of tensor rank
there. The polarisation class $\eta$ lies in $V_{+}\otimes V_{-}$ and is the
hermitian form of \Cref{prop:weilline}, hence nondegenerate: in dual bases
$\eta=\sum_{a=1}^{2n}v_{a}\wedge w_{a}$ with $v_{a}$ a basis of $V_{+}$ and
$w_{a}$ a basis of $V_{-}$. Squaring,
\[
  \eta^{2}
  =\sum_{a,b}(v_{a}\wedge w_{a})(v_{b}\wedge w_{b})
  =-2\sum_{a<b}(v_{a}v_{b})\otimes(w_{a}w_{b}),
\]
which is diagonal in the induced bases of $\bigwedge^{2}V_{+}$ and
$\bigwedge^{2}V_{-}$ and therefore has tensor rank exactly $\binom{2n}{2}$.
The class $\theta^{+}\theta^{-}$ is, up to sign,
$\theta^{+}\otimes\theta^{-}$, of tensor rank one. For $n\ge2$ one has
$\binom{2n}{2}\ge6>1$, so the two cannot be proportional.
\end{proof}

\begin{theorem}[The positivity obstruction]\label{thm:positivity}
Let $(A,\iota,\eta)$ be of $(K,-1,n)$-Weil type with $n\ge2$, and suppose that
the only divisor classes on which $\Kx$ acts through the norm character are
the multiples of $\eta$, that is
\begin{equation}\label{eq:noextranorm}
  \NS(A)_{\QQ}\cap\bigl(V_{+}\otimes V_{-}\bigr)=\QQ\,\eta .
\end{equation}
Let $E$ be a perfect complex on $A$ whose Chern character is effective in the
exponential basis, meaning
\begin{equation}\label{eq:effectivech}
  \ch(E)=\sum_{i=1}^{s}m_{i}\,e^{\lambda_{i}},
  \qquad m_{i}\in\QQ_{>0},\quad\lambda_{i}\in\NS(A)_{\QQ},
\end{equation}
and write $\lambda_{i}=e_{i}\eta+\theta_{i}$ with $\theta_{i}$ in the part $R$
of $H^{2}(A,\QQ)$ carrying the characters $\tau^{2}$ and $\bbar\tau^{\,2}$, and
$u_{i}\in K$ for the $\theta^{+}$-coordinate of $\theta_{i}$. If $\ch_{2}(E)$
is a flat Hodge class, then
\begin{equation}\label{eq:normsum}
  \sum_{i=1}^{s}m_{i}\,\Nm_{K/\QQ}(u_{i})=0 ,
\end{equation}
and since $\Nm_{K/\QQ}$ is positive definite and every $m_{i}$ is positive,
$u_{i}=0$ for every $i$. Hence every $\lambda_{i}$ is a rational multiple of
$\eta$, and $\ch_{k}(E)\in\QQ\,\eta^{k}$ for every $k$. In particular
$\ch_{n}(E)$ has no component on the Weil line.
\end{theorem}

\begin{proof}
From \eqref{eq:effectivech},
$\ch_{2}(E)=\tfrac12\sum_{i}m_{i}\lambda_{i}^{2}$, and hypothesis
\eqref{eq:noextranorm} is what makes the decomposition
$\lambda_{i}=e_{i}\eta+\theta_{i}$ available with $e_{i}$ a scalar. Decompose
$H^{4}(A,\CC)$ into $\Kx$-isotypic pieces and keep the one for $\Nm^{2}$,
namely $\bigwedge^{2}V_{+}\otimes\bigwedge^{2}V_{-}$. Writing
$\theta_{i}=u_{i}\theta^{+}+\bbar u_{i}\theta^{-}$ for a fixed choice of
generators $\theta^{\pm}$ of $\bigwedge^{2}V_{\pm}$, one has
\[
  \lambda_{i}^{2}
  =e_{i}^{2}\eta^{2}
   +2e_{i}\eta\theta_{i}
   +u_{i}^{2}(\theta^{+})^{2}+\bbar u_{i}^{2}(\theta^{-})^{2}
   +2\Nm_{K/\QQ}(u_{i})\,\theta^{+}\theta^{-},
\]
and of these five terms only the first and the last are in the $\Nm^{2}$
piece: $\eta\theta_{i}$ carries $\Nm\tau^{2}$ and $\Nm\bbar\tau^{\,2}$, and
$(\theta^{\pm})^{2}$ carry $\tau^{4}$ and $\bbar\tau^{\,4}$, all four distinct
from $\Nm^{2}$. So the $\Nm^{2}$-part of $\ch_{2}(E)$ is
\[
  \tfrac12\Bigl(\sum_{i}m_{i}e_{i}^{2}\Bigr)\eta^{2}
  +\Bigl(\sum_{i}m_{i}\Nm_{K/\QQ}(u_{i})\Bigr)\theta^{+}\theta^{-}.
\]

Now compare with the flat Hodge classes of degree four. For $n\ge3$ they are
$\QQ\eta^{2}$ by \Cref{thm:hdgcount}. For $n=2$ they are
$\QQ\eta^{2}\oplus\HW(A)$, and $\HW(A)$ spans the summands for $\tau^{4}$ and
$\bbar\tau^{\,4}$, which are not $\Nm^{2}$. In both cases the $\Nm^{2}$-part
of a flat Hodge class of degree four is a multiple of $\eta^{2}$. By
\Cref{lem:tensorrank} the class $\theta^{+}\theta^{-}$ is not such a multiple,
so its coefficient must vanish, which is \eqref{eq:normsum}.

Finally $\Nm_{K/\QQ}(g+h\sqrt{-d})=g^{2}+dh^{2}$ with $d>0$, so
$\Nm_{K/\QQ}(u)\ge0$ with equality only at $u=0$, and a sum of nonnegative
terms with positive weights vanishes only when every term does. Then
$\ch_{k}(E)=\frac{1}{k!}\bigl(\sum_{i}m_{i}e_{i}^{k}\bigr)\eta^{k}$, and
$\eta^{k}$ is not on the Weil line by \Cref{thm:hdgcount}.
\end{proof}

\begin{corollary}[Objects filtered by line bundles]\label{cor:nofiltered}
Let $A$ satisfy \eqref{eq:noextranorm}, which holds at every split member of
\Cref{thm:splitclosed} and at every quaternionic base point of
\Cref{thm:quatweil}. Then \Cref{thm:semiregclosure} cannot be applied at $A$
with any of the following:
\begin{enumerate}[label=\textup{(\alph*)},nosep]
\item a coherent sheaf admitting a filtration by line bundles, in particular
  any direct sum of line bundles;
\item a simple semi-homogeneous vector bundle in the sense of Mukai
  \cite{Muk78}, since such a bundle of rank $r$ has $\ch(E)=r\,e^{c_{1}(E)/r}$;
\item any iterated extension of the objects in (a) and (b).
\end{enumerate}
In each case the failure is not of the semiregularity hypothesis but of the
Chern character hypothesis of \Cref{prop:ktheory}: no such $E$ has a Chern
character which is a flat section of Hodge classes with a nonzero Weil
component.
\end{corollary}

\begin{proof}
First, \eqref{eq:noextranorm} holds at the points named. At a split member
$\NS(A)_{\QQ}$ has rank three with basis $\beta,\widehat\beta,\ell$ by
\Cref{prop:splitgeom}(iv), and of these $\eta$ spans the norm-character part
while $\gamma$ and $\ell$ span $R$. At a quaternionic point $\NS(A)_{\QQ}$ is
the space of Rosati-symmetric elements of $B=(-d,b)_{\QQ}$, which is spanned
by $1$, $j$ and $ij$; the first gives $\eta$ and the other two are semilinear
over $K$, hence lie in $R$. In both cases the norm-character part of
$\NS(A)_{\QQ}$ is exactly $\QQ\eta$.

Each of (a), (b) and (c) has a Chern character of the shape
\eqref{eq:effectivech} with positive multiplicities, by additivity of $\ch$
along a filtration and by Mukai's computation of the Chern character of a
simple semi-homogeneous bundle. The conclusion of \Cref{thm:positivity} then
contradicts the requirement that $\ch_{n}(E)$ have a nonzero component on the
Weil line.
\end{proof}

\begin{remark}[Scope of the theorem, and the objects of
Markman]\label{rem:markmanpositivity}
At a very general member of the family $\NS(A)_{\QQ}=\QQ\eta$, so
\Cref{thm:positivity} is true for the empty reason that there are no
$\theta_{i}$ to begin with. Its content is at the special members, which is
exactly where \Cref{thm:semiregclosure} needs its object to live and where the
base points of \Cref{thm:everyfamily} are. That is the opposite of the usual
situation in this paper, where the special members are the easy ones.

The theorem does not touch the objects of Markman \cite{Mar25a,Mar25b}.
Those are sheaves whose class in $K_{0}$ is not a positive combination of
exponentials of divisor classes, so \eqref{eq:effectivech} fails for them at
the first step; and they live at very general points, where, as just noted,
the statement is vacuous. What the theorem does is remove the class of
objects one would reach for first, and it removes it for a reason which is
visible before any cohomology is computed.
\end{remark}

\begin{remark}[Signed sums of line bundles]
\label{rem:whatitmustbe}
\Cref{thm:positivity} uses positivity of the multiplicities and nothing else,
so it says exactly where to go next: the object must be a complex whose
class in $K_{0}$ is not effective, that is one in which the line bundles
appear with multiplicities of both signs. Such objects exist. Writing
$E=F\oplus G[1]$ with $F$ and $G$ sums of line bundles, the class in $K_{0}$
is $[F]-[G]$ and \eqref{eq:normsum} becomes
$\sum_{L_{i}\subset F}\Nm(u_{i})=\sum_{L_{j}\subset G}\Nm(u_{j})$, which is
satisfiable; and \Cref{lem:splitsemireg} still applies verbatim, with the
shifts absorbing condition (ii).

The object of \Cref{thm:semiregclosure}, if it exists at a split member, is
not a sum of line bundles in any arrangement of signs and shifts:
\Cref{cor:noline} below proves this in every dimension. It must have
$\Ext^{2}$ that the trace does not see, that is, it must be indecomposable in
the derived category. This description of the remaining construction is
sharper than the one in \Cref{rem:tworemain}.
\end{remark}

\subsection{An explicit object and the Prouhet--Tarry--Escott problem}
\label{ssec:pte}

\Cref{prop:ktheory} produces the object of \Cref{thm:semiregclosure} by an
existence argument: the Chern character is an isomorphism onto the Chow ring
with rational coefficients, so a $K$-theory class with the right character
exists. That is enough for the theorem and tells one nothing about the object.
This subsection writes it down. The answer turns the question into a
Diophantine one, and the Diophantine problem has a classical name.

\begin{setupenv}[Objects supported on $R$]\label{setup:pte}
Let $A$ be of $(K,-1,n)$-Weil type, $n\ge2$, satisfying
\eqref{eq:noextranorm}, and keep $\theta^{\pm}$ from
\Cref{not:splitcoords}. For $u\in K$ with $\theta(u)=u\theta^{+}+\bbar
u\,\theta^{-}\in\NS(A)_{\QQ}$, and for multiplicities $m_{i}\in\ZZ$, consider
\[
  v=\sum_{i=1}^{s}m_{i}\,e^{\theta(u_{i})},
  \qquad
  T_{q,r}=\sum_{i=1}^{s}m_{i}\,u_{i}^{q}\,\bbar u_{i}^{\,r}
  \quad (q,r\ge0).
\]
\end{setupenv}

\begin{proposition}[The Chern character in terms of the $u_{i}$]
\label{prop:chintermsofu}
In \Cref{setup:pte},
\begin{equation}\label{eq:chexpansion}
  v_{k}
  =\sum_{\substack{q+r=k\\ 0\le q,r\le n}}
    \frac{T_{q,r}}{q!\,r!}\,(\theta^{+})^{q}(\theta^{-})^{r},
\end{equation}
because $(\theta^{\pm})^{q}=0$ for $q>n$. Consequently $v$ is a flat section
of Hodge classes with a nonzero Weil component if
\begin{equation}\label{eq:pteconditions}
  T_{q,r}=0\ \ \text{for all } 0\le q,r\le n \text{ with }
  (q,r)\notin\{(0,0),(n,0),(0,n)\},
  \qquad T_{n,0}\ne0,
\end{equation}
and then $v_{0}=T_{0,0}=\sum_{i}m_{i}$, $v_{k}=0$ for $0<k<2n$ with $k\ne n$,
and
\[
  v_{n}=\frac{1}{n!}\Bigl(T_{n,0}(\theta^{+})^{n}
        +\bbar{T}_{n,0}(\theta^{-})^{n}\Bigr),
\]
a nonzero rational multiple of a generator of the Weil line, by
\Cref{thm:splitclosed}.
\end{proposition}

\begin{proof}
Expand $\theta(u_{i})^{k}$ by the binomial theorem, divide by $k!$ and collect;
the vanishing of $(\theta^{\pm})^{q}$ for $q>n$ is because
$\theta^{\pm}\in\bigwedge^{2}V_{\pm}$ and $\dim_{\CC}V_{\pm}=2n$. Under
\eqref{eq:pteconditions} the only surviving terms of \eqref{eq:chexpansion}
are the ones listed, and $(\theta^{+})^{n}$ spans $\bigwedge^{2n}V_{+}$, which
is the Weil line over $\CC$.
\end{proof}

So the search for the object of \Cref{thm:semiregclosure} is the search for a
finitely supported integer-valued function on $\cO_{K}$ that annihilates every
monomial $u^{q}\bbar u^{\,r}$ with $q,r\le n$ except three of them. Two things
follow at once, and they pull in opposite directions.

\begin{lemma}[Equal norms force the Prouhet--Tarry--Escott condition]
\label{lem:pte}
Suppose all the $u_{i}$ in \Cref{setup:pte} have the same norm
$N=\Nm_{K/\QQ}(u_{i})>0$. Then \eqref{eq:pteconditions} holds if and only if
\begin{equation}\label{eq:ptepure}
  \sum_{i}m_{i}u_{i}^{\,k}=0\ \ \text{for } k=0,1,\dots,n-1,
  \qquad
  \sum_{i}m_{i}u_{i}^{\,n}\ne0 .
\end{equation}
With $m_{i}=\pm1$ this says that the $u_{i}$ split into two multisets on the
circle $\Nm=N$ whose power sums agree through degree $n-1$ and differ at
degree $n$: a Prouhet--Tarry--Escott pair of degree $n-1$, constrained to a
norm circle of $\cO_{K}$.
\end{lemma}

\begin{proof}
For $q\ge r$ one has $u^{q}\bbar u^{\,r}=N^{r}u^{q-r}$, so
$T_{q,r}=N^{r}\sum_{i}m_{i}u_{i}^{\,q-r}$, and the pairs $(q,r)$ with
$q,r\le n$ realise every difference $k=q-r$ from $0$ to $n$. The excluded
pairs $(0,0)$ and $(n,0)$ are the ones with $r=0$ and $k=0$ or $k=n$; but
$k=0$ also occurs at $(1,1),\dots,(n,n)$, which are not excluded, so
$\sum_{i}m_{i}=0$ is still required. The remaining differences
$k=1,\dots,n-1$ each occur at some non-excluded pair, and $k=n$ occurs only at
$(n,0)$.
\end{proof}

\begin{theorem}[Equal norms cannot be semiregular]\label{thm:equalnorms}
Let $E=\bigoplus_{i}L_{i}[p_{i}]$ realise \eqref{eq:ptepure} with all norms
equal. Then $\ch_{0}(E)=0$, and the semiregularity map $\sigma_{E}$ has a
kernel of dimension at least $n^{2}$.
\end{theorem}

\begin{proof}
By \Cref{lem:pte} the case $k=0$ of \eqref{eq:ptepure} gives
$\sum_{i}m_{i}=0$, which is $\ch_{0}(E)$. For the kernel, take
$\zeta_{i}=m_{i}\zeta$ for a single $\zeta\in H^{2}(\cO_{A})$; then by
\Cref{lem:splitsemireg}(iii)
\[
  \Phi\bigl((m_{i}\zeta)_{i}\bigr)
  =\Bigl(\sum_{i}m_{i}e^{\lambda_{i}}\Bigr)\zeta=\ch(E)\,\zeta ,
\]
which by \Cref{prop:chintermsofu} is a multiple of $\omega\,\zeta$ with
$\omega$ a generator of the Weil line. Now $\HW(A)_{\CC}$ is spanned by
$\bigwedge^{2n}V_{+}$ and $\bigwedge^{2n}V_{-}$, and the Weil type condition
says $\dim V_{\pm}^{0,1}=n$; writing
$H^{2}(\cO_{A})=\bigwedge^{2}(V_{+}^{0,1}\oplus V_{-}^{0,1})$, the product of
a generator of $\bigwedge^{2n}V_{+}$ with $\zeta$ kills the components of
$\zeta$ in $\bigwedge^{2}V_{+}^{0,1}$ and in
$V_{+}^{0,1}\otimes V_{-}^{0,1}$, and likewise with the roles exchanged. So
$\omega\,\zeta=0$ for every $\zeta$ in
$V_{+}^{0,1}\otimes V_{-}^{0,1}$, a space of dimension $n^{2}$, and all of
those lie in the kernel of $\Phi$.
\end{proof}

\begin{corollary}[At least $n+1$ norms]\label{cor:nplusone}
If $E$ as in \Cref{setup:pte} satisfies \eqref{eq:pteconditions} and
$\ch_{0}(E)\ne0$, then the $u_{i}$ take at least $n+1$ distinct norms.
\end{corollary}

\begin{proof}
Let $N_{1},\dots,N_{\nu}$ be the distinct norms and $M_{j}$ the sum of the
$m_{i}$ over the level $\Nm(u_{i})=N_{j}$. The conditions $T_{r,r}=0$ for
$r=1,\dots,n$ read $\sum_{j}M_{j}N_{j}^{\,r}=0$ for $r=1,\dots,n$. The matrix
$(N_{j}^{\,r})_{1\le r\le n,\,1\le j\le\nu}$ has rank $\min(n,\nu)$ because the
$N_{j}$ are distinct and nonzero, so $\nu\le n$ forces every $M_{j}=0$ and
hence $\ch_{0}(E)=\sum_{j}M_{j}=0$.
\end{proof}

\subsubsection*{The smallest solution}

At $n=3$ and $K=\QQ(\sqrt{-3})$ the Prouhet--Tarry--Escott condition of
\Cref{lem:pte} has a particularly simple solution: the norm circle
$\Nm=1$ carries the six roots of unity of $K$, and the two cosets of the cube
roots of $-1$ and of $+1$ are a Prouhet--Tarry--Escott pair of degree
two. Explicitly, with $\zeta_{6}$ a primitive sixth root of unity, the
cube roots of $-1$ are $\zeta_{6},\zeta_{6}^{3},\zeta_{6}^{5}$ and those of
$+1$ are $1,\zeta_{6}^{2},\zeta_{6}^{4}$; each triple sums to zero and has
zero sum of squares,
while the sums of cubes are $-3$ and $+3$.

\begin{proposition}[An explicit object at $n=3$]\label{prop:explicitobject}
Let $K=\QQ(\sqrt{-3})$, $n=3$, and let $A_{0}=X\times\widehat X$ be a split
member. In the basis $\beta,\widehat\beta,\ell$ of $\NS(A_{0})_{\QQ}$ put
\[
  \lambda_{1}=(3,-1,3),\quad
  \lambda_{2}=(-6,2,0),\quad
  \lambda_{3}=(3,-1,-3),
\]
and $\lambda_{3+i}=-\lambda_{i}$ for $i=1,2,3$, with $m_{i}=+1$ for
$i\le3$ and $m_{i}=-1$ for $i\ge4$. Then the six classes are distinct and
integral, the perfect complex
\[
  E=\bigl(L_{1}\oplus L_{2}\oplus L_{3}\bigr)\ \oplus\
    \bigl(L_{4}\oplus L_{5}\oplus L_{6}\bigr)[3]
\]
has
\[
  \ch_{k}(E)=0\ \ (k\ne3),
  \qquad
  \ch_{3}(E)=12\,\omega_{1},
\]
and every difference $\lambda_{j}-\lambda_{i}$ with $i\ne j$ is nondegenerate
of index exactly $3$, so that conditions (i) and (ii) of
\Cref{lem:splitsemireg} hold for this shift. Condition (iii) fails: the map
$\Phi$ has rank $75$ instead of $90$, with kernel of dimension
$15=\dim H^{2}(\cO_{A_{0}})$.
\end{proposition}

\begin{proof}
With $d=3$, the class $\theta(u)=2g\gamma+6h\ell$ of $u=g+h\sqrt{-3}$ has
coordinates $(6g,-2g,6h)$, so $\lambda_{1},\lambda_{2},\lambda_{3}$ are the
images of the cube roots $\zeta_{6},\zeta_{6}^{3},\zeta_{6}^{5}$ of $-1$ and
$\lambda_{4},\lambda_{5},\lambda_{6}$ those of the cube roots of $+1$; the six
classes are distinct and integral. All six have norm $1$, and
$\sum_{i}m_{i}u_{i}^{k}$ is $0$ for $k=0,1,2$ and $-3-3=-6$ for $k=3$, so
\eqref{eq:ptepure} holds at $n=3$. By \Cref{lem:pte} and
\Cref{prop:chintermsofu}, $\ch_{k}(E)=0$ for $k\ne3$ and
$\ch_{3}(E)=-\bigl((\theta^{+})^{3}+(\theta^{-})^{3}\bigr)$, and by
\Cref{ex:splitsmall} at $n=3$,
$\omega_{1}\pm\sqrt{-3}\,\omega_{2}=-\frac16(\theta^{\mp})^{3}$; so
$\ch_{3}(E)=12\,\omega_{1}$.

Every difference is $\lambda_{j}-\lambda_{i}=\theta(w)$ with
$w=u_{j}-u_{i}\ne0$, and $\theta(w)$ is nondegenerate of index $n$ for every
$w\in K^{\times}$. Indeed $\theta(1)=2\gamma$ and $\theta(\sqrt{-d})=2d\ell$ are of
type $(1,1)$, hence so are $\theta^{\pm}$, so
$\theta^{+}\in V_{+}^{1,0}\otimes V_{+}^{0,1}$, and $V_{+}^{0,1}$ is the
conjugate of $V_{-}^{1,0}$. In the splitting of the tangent space dual to
$H^{1,0}=V_{+}^{1,0}\oplus V_{-}^{1,0}$, the hermitian form of $\theta(w)$
therefore has the block form
$\bigl(\begin{smallmatrix}0&wP\\ \bar wP^{*}&0\end{smallmatrix}\bigr)$, with $P$
the block of $2\gamma$, which is invertible because $\gamma$ is nondegenerate.
Its eigenvalues are plus and minus the singular values of $wP$, so it has
signature $(n,n)$, and a line bundle with this class has cohomology only in
degree $n$ \cite[Chapter~3]{BL04}. For the shift, condition (ii) asks that
$2+p_{i}-p_{j}\ne3$ whenever $i\ne j$, that is $p_{i}-p_{j}\ne1$; within each
of the two blocks the difference is $0$, and between them it is $\pm3$.

For the rank, the $u_{i}$ are the six sixth roots of unity, each once, and
$\bbar u_{i}=u_{i}^{-1}$. Absorb the signs into the $\zeta_{i}$ as in
\Cref{lem:splitsemireg}(iii) and expand
$e^{\lambda_{i}}=e^{u_{i}\theta^{+}}e^{\bbar u_{i}\theta^{-}}$:
\[
  \Phi(\zeta)=\sum_{a,b=0}^{3}\frac{(\theta^{+})^{a}(\theta^{-})^{b}}{a!\,b!}\,Z_{a-b},
  \qquad Z_{k}=\sum_{i}u_{i}^{k}\zeta_{i},
\]
with $k$ read modulo $6$. The $Z_{0},\dots,Z_{5}$ determine the $\zeta_{i}$,
by a Vandermonde matrix. Split
$H^{0,2}(A_{0})=\bigwedge^{2}V_{+}^{0,1}\oplus(V_{+}^{0,1}\otimes V_{-}^{0,1})
\oplus\bigwedge^{2}V_{-}^{0,1}$, write $Z_{k}=Z_{k}^{+}+Z_{k}^{0}+Z_{k}^{-}$
accordingly, and grade $H^{*}(A_{0})=\bigwedge^{*}V_{+}\otimes\bigwedge^{*}V_{-}$
by the Hodge types of the two factors. The term
$(\theta^{+})^{a}(\theta^{-})^{b}Z_{k}^{0}$ has types $(a,a+1)$ and $(b,b+1)$,
the term with $Z_{k}^{+}$ has types $(a,a+2)$ and $(b,b)$, and the term with
$Z_{k}^{-}$ has types $(a,a)$ and $(b,b+2)$. These are different for
different $(a,b)$ and different pieces, so every term vanishes separately.
Since $\theta^{\pm}$ pairs $V_{\pm}^{1,0}$ nondegenerately with $V_{\pm}^{0,1}$,
both of dimension $3$, multiplication by $(\theta^{\pm})^{a}$ is injective on
$\bigwedge^{j}V_{\pm}^{0,1}$ when $a+j\le3$ and zero otherwise; in bases
with $\theta^{+}=\sum_{i}e_{i}f_{i}$ it sends $f_{J}$ to a signed sum of the
$e_{I}f_{I\cup J}$ with $I\cap J=\emptyset$, $|I|=a$, and different $J$ give
disjoint supports. Hence $Z_{k}^{0}=0$ unless no $a,b\le2$ have $a-b\equiv k$,
which leaves $k=3$; $Z_{k}^{+}=0$ unless no $a\le1$, $b\le3$ have
$a-b\equiv k$, which leaves $k=2$; and likewise $Z_{k}^{-}=0$ unless $k=4$. So
$\ker\Phi$ is $V_{+}^{0,1}\otimes V_{-}^{0,1}$ at $k=3$ together with
$\bigwedge^{2}V_{+}^{0,1}$ at $k=2$ and $\bigwedge^{2}V_{-}^{0,1}$ at $k=4$, of
dimension $9+3+3=15=\dim H^{2}(\cO_{A_{0}})$, and $\Phi$ has rank
$90-15=75$. The first summand is the kernel of \Cref{thm:equalnorms}.
\end{proof}

\begin{remark}[What the construction settles]\label{rem:pteoutcome}
Three things are settled. The object of \Cref{prop:ktheory} is not an
abstraction: at $n=3$ over $\QQ(\sqrt{-3})$ it is six line bundles and a
shift, and its Chern character is exactly $12$ times a Weil class with every
other graded piece zero. The construction is completely general, and
\Cref{lem:pte} identifies the constraint as a Prouhet--Tarry--Escott problem
on a norm circle, a question of a kind that can be attacked directly. And
\Cref{cor:nplusone} says what any semiregular solution must look like: the
$u_{i}$ must take at least $n+1$ distinct norms, because otherwise the rank is
zero and the argument of \Cref{thm:equalnorms} produces a kernel of
dimension $n^{2}$.

This leaves the existence of a solution with $n+1$ norms. The level sums are
then forced up to scale (\Cref{thm:pterange}), and \Cref{cor:noline} shows
that no solution could serve, because a line bundle summand off the
polarisation line is obstructed along the family whatever the norms.
\end{remark}

\subsubsection*{The obstruction classes of the explicit object}

\Cref{cor:kernelfamily} and \Cref{rem:whatmeasures} say that the kernel of
$\sigma_{E}$ contains a copy of the tangent space of the Weil family, and
that this is a statement about where the semiregularity map is blind, not a
statement that the object is obstructed. For the explicit object the
obstruction classes can be written down, and they are not zero. The
computation is elementary and we record it in general.

Recall the conventions of \Cref{ssec:tangent}: a first-order deformation of
$A$ is $v\in\Hom(H^{1,0},H^{0,1})=H^{1}(A,T_{A})$, extended by zero on
$H^{0,1}$ and acting on $H^{*}(A,\CC)$ as a derivation $v\lrcorner$; a class
$\xi$ of type $(p,p)$ stays of type $(p,p)$ to first order along $v$ exactly
when $v\lrcorner\,\xi=0$. For a line bundle $L$ with $\lambda=c_{1}(L)$ the
Atiyah class is $\lambda\in H^{1}(\Omega^{1}_{A})$, its contraction with $v$
is $v\lrcorner\,\lambda\in H^{2}(\cO_{A})=H^{0,2}(A)$, and this is the
obstruction to deforming $L$ along $v$; a shift does not change it. For
$E=\bigoplus_{i}L_{i}[p_{i}]$ the Atiyah class is diagonal with entries
$\lambda_{i}$, as in the proof of \Cref{lem:splitsemireg}, so the
obstruction to deforming $E$ along $v$ is the diagonal element
\begin{equation}\label{eq:obdef}
  \ob_{E}(v)=\bigl(v\lrcorner\,\lambda_{i}\bigr)_{i}
  \in\bigoplus_{i}H^{0,2}(A)\subset\Ext^{2}(E,E).
\end{equation}

\begin{proposition}[The obstruction map of a split object]
\label{prop:splitobstruction}
Let $A_{0}=X\times\widehat X$ be a split member with $\beta$ principal and
$X$ general, so that $\NS(A_{0})_{\QQ}=\langle\beta,\widehat\beta,\ell\rangle$
by \Cref{prop:splitgeom}(iv). Let $T=T_{[A_{0}]}\cD_{n,n}$ be the space of
polarised first-order deformations commuting with $K$, of dimension $n^{2}$
by \Cref{prop:firstorder}, and put
\[
  T_{\mathrm{spl}}
  =\bigl\{\,v\in\Hom(H^{1,0},H^{0,1})\;:\;
    v\lrcorner\,\beta=v\lrcorner\,\widehat\beta=v\lrcorner\,\ell=0\,\bigr\}.
\]
Let $E=\bigoplus_{i=1}^{s}L_{i}[p_{i}]$ with $\lambda_{i}=c_{1}(L_{i})$ and
let $\ob_{E}\colon T\to\bigoplus_{i}H^{0,2}(A_{0})$ be the map
\eqref{eq:obdef} restricted to $T$.
\begin{enumerate}[label=\textup{(\roman*)},nosep]
\item $T_{\mathrm{spl}}$ is the tangent space to the split locus: it has
  dimension $n(n+1)/2$, it is contained in $T$, and it consists of the
  deformations of $X\times\widehat X$ induced by the polarised deformations
  of $(X,\beta)$.
\item If the classes $\lambda_{i}$ together with $\eta$ span
  $\NS(A_{0})_{\QQ}$, then $\ker\ob_{E}=T_{\mathrm{spl}}$ and
  $\operatorname{rank}\ob_{E}=n(n-1)/2$, the codimension of the split locus
  in $\cD_{n,n}$.
\item If $\ch_{k}(E)\in\QQ\eta^{k}$ for $k\ne n$ and
  $\ch_{n}(E)\in\QQ\eta^{n}\oplus\HW(A_{0})$, then
  $\sigma_{E}\circ\ob_{E}=0$ on $T$: for every $v\in T$,
  \begin{equation}\label{eq:sigmaob}
    \sigma_{E}\bigl(\ob_{E}(v)\bigr)=v\lrcorner\,\ch(E)=0 .
  \end{equation}
\item For the object of \Cref{prop:explicitobject}, at $n=3$ and
  $K=\QQ(\sqrt{-3})$: $\ker\ob_{E}=T_{\mathrm{spl}}$ has dimension $6$ and
  $\ob_{E}$ has rank $3$; the image of $\ob_{E}$ lies in $\ker\sigma_{E}$;
  the subspace $\Delta=\{z\cdot\id_{E}:z\in\Ann_{H^{0,2}}(\omega_{1})\}$ of
  \Cref{cor:kernelfamily} has dimension $9$ and meets the image of $\ob_{E}$
  only in zero, so that $\Delta\oplus\image\ob_{E}$ is a twelve-dimensional
  subspace of the fifteen-dimensional kernel of $\sigma_{E}$; and
  $\ob_{E}(v)$ has components $v\lrcorner\,\lambda_{i}$ on $L_{i}$ and
  $-v\lrcorner\,\lambda_{i}$ on $L_{i}^{-1}[3]$, so it lies on the line of
  vectors of the shape $(z,\dots,z)$ only when it vanishes. In particular
  $E$ has a nonzero first-order obstruction along every $v\in T$ not tangent
  to the split locus, and no element $z\cdot\id_{E}$ of $\Ext^{2}(E,E)$
  cancels it.
\end{enumerate}
\end{proposition}

\begin{proof}
(i) Let $P=H^{1,0}(X)\subset H^{1}(X,\CC)=U^{\vee}_{\CC}$ with basis
$p_{1},\dots,p_{n}$, so that $p_{1},\dots,p_{n},\bbar p_{1},\dots,\bbar p_{n}$
is a basis of $U^{\vee}_{\CC}$, and let $q_{a},\bbar q_{a}$ be the dual basis
of $U_{\CC}=H^{1}(\widehat X,\CC)$; the dual of a conjugate basis is the
conjugate of the dual basis, so the notation is consistent. The Poincar\'e
class is the identity element of $U^{\vee}\otimes U$,
\[
  \ell=\sum_{a=1}^{n}\bigl(p_{a}\wedge q_{a}+\bbar p_{a}\wedge\bbar q_{a}\bigr),
\]
and it is of type $(1,1)$, which forces $H^{1,0}(\widehat X)=\langle\bbar
q_{a}\rangle$ and $H^{0,1}(\widehat X)=\langle q_{a}\rangle$: the Hodge
structure of $\widehat X$ is the dual one. Since $\beta$ is of type $(1,1)$
and nondegenerate, $\beta=\sum_{a,b}h_{ab}\,p_{a}\wedge\bbar p_{b}$ with
$h\in\GL_{n}(\CC)$; and since $\beta$ is principal with Gram matrix $B$ in a
symplectic basis of $U$ and $\widehat\beta$ has Gram matrix $-B^{-1}=B$ in
the dual basis of $U^{\vee}$, the Gram matrix of $\widehat\beta$ in the
basis $p_{a},\bbar p_{a}$ is minus the inverse of the Gram matrix
$\begin{psmallmatrix}0&h\\-h^{\top}&0\end{psmallmatrix}$ of $\beta$ in the
basis $q_{a},\bbar q_{a}$, which gives
$\widehat\beta=-\sum_{a,b}(h^{-1})_{ab}\,\bbar q_{a}\wedge q_{b}$.

Now $H^{1,0}(A_{0})=\langle p_{a}\rangle\oplus\langle\bbar q_{a}\rangle$ and
$H^{0,1}(A_{0})=\langle\bbar p_{a}\rangle\oplus\langle q_{a}\rangle$, so a
deformation $v$ is four $n\times n$ matrices,
\[
  v(p_{a})=\sum_{c}\bigl(A_{ca}\,\bbar p_{c}+C_{ca}\,q_{c}\bigr),
  \qquad
  v(\bbar q_{a})=\sum_{c}\bigl(B_{ca}\,\bbar p_{c}+D_{ca}\,q_{c}\bigr),
\]
and $v$ kills $\bbar p_{a}$ and $q_{a}$. Contracting,
\[
  v\lrcorner\,\ell
  =\sum_{a,c}\Bigl(A_{ca}\,\bbar p_{c}\wedge q_{a}+C_{ca}\,q_{c}\wedge q_{a}
   +B_{ca}\,\bbar p_{a}\wedge\bbar p_{c}+D_{ca}\,\bbar p_{a}\wedge q_{c}\Bigr),
\]
and reading off the three kinds of monomials, $v\lrcorner\,\ell=0$ if and only
if $D=-A^{\top}$ and $B$, $C$ are symmetric. Next
\[
  v\lrcorner\,\beta
  =\sum_{a,b,c}h_{ab}\Bigl(A_{ca}\,\bbar p_{c}\wedge\bbar p_{b}
   +C_{ca}\,q_{c}\wedge\bbar p_{b}\Bigr),
\]
so $v\lrcorner\,\beta=0$ if and only if $Ch=0$, that is $C=0$, and $Ah$ is
symmetric. Likewise $v\lrcorner\,\widehat\beta=0$ if and only if $B=0$ and
$Dh^{-1}$ is symmetric. Given $D=-A^{\top}$, the last condition reads
$A^{\top}h^{-1}=h^{-\top}A$, which after multiplication by $h^{\top}$ on the
left and by $h$ on the right is $h^{\top}A^{\top}=Ah$, the symmetry of $Ah$
again. Hence
\[
  T_{\mathrm{spl}}=\bigl\{(A,B,C,D)=(A,0,0,-A^{\top})\;:\;Ah\ \text{symmetric}\bigr\},
\]
of dimension $n(n+1)/2$. It is the tangent space to the split locus: the
same computation on $X$ alone shows that the polarised first-order
deformations of $(X,\beta)$ are the matrices $A$ with $Ah$ symmetric, the
space $T_{[X]}\cA_{n}$ of dimension $n(n+1)/2$; the deformation of
$\widehat X$ induced by $A$ is the dual one, $-A^{\top}$; and the map
$A\mapsto(A,0,0,-A^{\top})$ is injective. Finally $T_{\mathrm{spl}}\subseteq
T$: for $v\in T_{\mathrm{spl}}$, $v\lrcorner\,\eta=0$ because
$\eta\in\langle\beta,\widehat\beta\rangle$, so $v$ is polarised, and
$v\lrcorner\,\omega_{1}=v\lrcorner\,\omega_{2}=0$ because $\omega_{1}$ and
$\omega_{2}$ are polynomials in $\beta,\widehat\beta,\ell$ by
\Cref{thm:splitclosed} and $v\lrcorner$ is a derivation; by
\Cref{prop:firstorder} such a $v$ commutes with $K$.

(ii) Every $v\in T$ has $v\lrcorner\,\eta=0$. If the $\lambda_{i}$ and
$\eta$ span $\NS(A_{0})_{\QQ}=\langle\beta,\widehat\beta,\ell\rangle$, then
$v\lrcorner\,\lambda_{i}=0$ for all $i$ is equivalent, for $v\in T$, to
$v\lrcorner\,\beta=v\lrcorner\,\widehat\beta=v\lrcorner\,\ell=0$, so
$\ker\ob_{E}=T\cap T_{\mathrm{spl}}=T_{\mathrm{spl}}$ by (i), and the rank
is $n^{2}-n(n+1)/2=n(n-1)/2$, which is the codimension of
\Cref{prop:splitgeom}(iii).

(iii) By \eqref{eq:splitsigma}, $\sigma_{E}(\zeta)=\sum_{i}(-1)^{p_{i}}
\zeta_{ii}\,e^{\lambda_{i}}$ for a diagonal $\zeta$. Since $v\lrcorner$ is a
derivation and $v\lrcorner\,\lambda_{i}$ has even degree,
$v\lrcorner\,\lambda_{i}^{\,q}=q\,\lambda_{i}^{\,q-1}(v\lrcorner\,\lambda_{i})$
and so $v\lrcorner\,e^{\lambda_{i}}=(v\lrcorner\,\lambda_{i})\,e^{\lambda_{i}}$.
Therefore
\[
  \sigma_{E}\bigl(\ob_{E}(v)\bigr)
  =\sum_{i}(-1)^{p_{i}}(v\lrcorner\,\lambda_{i})\,e^{\lambda_{i}}
  =v\lrcorner\sum_{i}(-1)^{p_{i}}e^{\lambda_{i}}
  =v\lrcorner\,\ch(E).
\]
For $v\in T$ one has $v\lrcorner\,\eta^{k}=k\eta^{k-1}(v\lrcorner\,\eta)=0$,
and $v\lrcorner\,\omega=0$ for every $\omega\in\HW(A_{0})$ by
\Cref{prop:firstorder}; under the hypothesis on $\ch(E)$ this gives
\eqref{eq:sigmaob}.

(iv) The three classes $\lambda_{1}=\gamma+3\ell$, $\lambda_{2}=-2\gamma$,
$\lambda_{3}=\gamma-3\ell$ span $\langle\gamma,\ell\rangle$, and
$\langle\eta,\gamma\rangle=\langle\beta,\widehat\beta\rangle$ by the
determinant in the proof of \Cref{lem:sixmonomials}, so (ii) applies and
gives the kernel and the rank; \Cref{prop:explicitobject} gives the
hypothesis of (iii), so the image lies in $\ker\sigma_{E}$. The subspace
$\Delta$ has dimension $\dim\Ann_{H^{0,2}}(\omega_{1})=n^{2}=9$ by
\Cref{lem:annihilator}. The components of $\ob_{E}(v)$ on $L_{i}$ and on
$L_{i}^{-1}[3]$ are $v\lrcorner\,\lambda_{i}$ and
$v\lrcorner\,(-\lambda_{i})=-v\lrcorner\,\lambda_{i}$, so an element of the
image of the shape $(z,\dots,z)$ has $z=-z$, hence is zero; this is
$\Delta\cap\image\ob_{E}=0$, and with $\dim\ker\sigma_{E}=15$ from
\Cref{prop:explicitobject} the dimension count follows.
\end{proof}

\begin{remark}[Meaning of the obstruction classes]\label{rem:obstructionmeaning}
\Cref{prop:splitobstruction} sharpens \Cref{rem:whatmeasures} in one
respect and exposes a limit of the criterion in another. First, for the
explicit object the obstruction classes along the family are nonzero, and
they are nonzero exactly in the $n(n-1)/2$ directions normal to the split
locus, which are the directions in which the Weil class has to be
transported. Second, those classes lie in the kernel of $\sigma_{E}$ by
\eqref{eq:sigmaob}, and \eqref{eq:sigmaob} holds for every object whose
Chern character is a flat section of Hodge classes along the family, so the
semiregularity map cannot detect the obstruction of any such object. The
obstruction is therefore nonzero, and the criterion cannot see it. For
the explicit object the
situation is the worst possible: the twelve dimensions
$\Delta\oplus\image\ob_{E}$ of the kernel are accounted for by the two
mechanisms of \Cref{cor:kernelfamily} and \eqref{eq:obdef}, the remaining
three lie on the even side of \Cref{cor:pteobstructed} as well, and none of
the fifteen is an artefact. The classes of the shape $z\cdot\id_{E}$ are the
ones a change of the twisting class would contribute, one and the same $z$
to every rank one summand; the last assertion of (iv) says that no such
change removes the obstruction. So the explicit object does not move off the
split locus to first order, either as a complex or as a twisted complex, and
the question of \Cref{q:residual} for it is not one of verifying a
hypothesis but of replacing the object.

The positive content is a recipe. By (ii), a split object whose classes
span $\NS(A_{0})_{\QQ}$ with $\eta$ has obstruction map of rank
$n(n-1)/2$ against the family, and the spanning hypothesis is more than is
needed: a single class $\theta\in\NS(A_{0})_{\QQ}$ not proportional to
$\eta$ already has $\{v\in T:v\lrcorner\,\theta=0\}=T_{\mathrm{spl}}$.
Indeed, every $v\in T$ kills $\eta$ and commutes with $K$, so it is a pair
$(v_{+},v_{-})$ with $v_{\pm}\in\Hom(V_{\pm}^{1,0},V_{\pm}^{0,1})$. Since
$\theta^{\pm}$ is of type $(1,1)$ in $\bigwedge^{2}V_{\pm}$, the $R_{0}$-part
$u\theta^{+}+\bbar u\,\theta^{-}$ of $\theta$, with $u\ne0$, gives
$v\lrcorner\,\theta=u\,v_{+}\lrcorner\,\theta^{+}+\bbar u\,v_{-}\lrcorner\,\theta^{-}$,
whose two terms lie in the different summands $\bigwedge^{2}V_{+}^{0,1}$ and
$\bigwedge^{2}V_{-}^{0,1}$ of $H^{0,2}(A_{0})$. So the kernel of
$v\mapsto v\lrcorner\,\theta$ on $T$ does not depend on $u\ne0$; it is
therefore the common kernel for $\gamma$ and $\ell$, the cases $u=1$ and
$u=\sqrt{-d}$, which is $T_{\mathrm{spl}}$ by (ii). So every split object built from line bundles
other than powers of $\eta$ is obstructed in exactly the normal directions
to the split locus, and \Cref{thm:positivity} says that an object built
from powers of $\eta$ carries no Weil class. What survives is the statement
of \Cref{rem:familyuniform}: the object must have positive rank.
\end{remark}

\subsubsection*{The evaluation map of the explicit object}

\Cref{thm:weakfails} places the copy $\Delta$ of the tangent space inside the
image of the evaluation map. For a split object the whole of
$\operatorname{ev}_{E}$ can be written down, and for the explicit object the
entire kernel of $\sigma_{E}$ turns out to lie in its image. We identify
$HH^{2}(A_{0})$ with
\begin{equation}\label{eq:HT2}
  HT^{2}(A_{0})
  =H^{0,2}(A_{0})\;\oplus\;\Hom\bigl(H^{1,0},H^{0,1}\bigr)\;\oplus\;
   \textstyle\bigwedge^{2}\bigl(H^{1,0}\bigr)^{\vee}
\end{equation}
by the Hochschild--Kostant--Rosenberg isomorphism, the three summands being
$H^{2}(\cO)$, $H^{1}(T)$ and $H^{0}(\bigwedge^{2}T)$ for the trivial tangent
bundle. An element $x=(z,v,\pi)$ acts on $H^{*}(A_{0},\CC)$ by $z\wedge$, by
the derivation $v\lrcorner$ of \Cref{ssec:tangent}, and by the double
contraction $\pi\lrcorner$, which removes two $(1,0)$ slots; we write
$x\lrcorner\,\xi$ for the sum of the three.

The double contraction is a second order operator. For a class $\lambda$ of
type $(1,1)$, which has a single $(1,0)$ slot, $\pi\lrcorner\,\lambda=0$ and
$\pi\lrcorner\,\lambda^{k}=\binom{k}{2}\lambda^{k-2}\,(\pi\lrcorner\,\lambda^{2})$,
because $\pi\lrcorner$ is a composite of two interior products, each of which
is an antiderivation, and the only nonzero terms are those in which the two
interior products hit two different factors. Summing over $k$ with the
factorials gives
$\pi\lrcorner\,e^{\lambda}=\tfrac12(\pi\lrcorner\,\lambda^{2})\,e^{\lambda}$,
exactly parallel to $v\lrcorner\,e^{\lambda}=(v\lrcorner\,\lambda)\,e^{\lambda}$
for the first order operator $v\lrcorner$. This is the whole content of the
compatibility \eqref{eq:sigmaev} below: the semiregularity map of a line
bundle is multiplication by $e^{\lambda}$, and every summand of $HT^{2}$ acts
on $e^{\lambda}$ as multiplication by a class in $H^{0,2}$ which is precisely
the value of $\operatorname{ev}_{L}$. The same identity, for arbitrary $E$, is
the commutative square of Buchweitz and Flenner \cite[Corollary 4.3]{BF03} and
its Hochschild form in \cite[Lemma 11.3]{Mar25b}; on an abelian variety the
Todd class is trivial and no correction to the HKR map is needed
\cite{Cal05,BF08}.

\begin{proposition}[The evaluation map of a split object]\label{prop:evaluation}
Let $A_{0}$ be a split member with $X$ general and let
$E=\bigoplus_{i=1}^{s}L_{i}[p_{i}]$ with $\lambda_{i}=c_{1}(L_{i})$ and all
$\Ext^{2}(L_{i}[p_{i}],L_{j}[p_{j}])$, $i\ne j$, zero, so that
$\Ext^{2}(E,E)=\bigoplus_{i}H^{0,2}(A_{0})$.
\begin{enumerate}[label=\textup{(\roman*)},nosep]
\item Under \eqref{eq:HT2}, $\operatorname{ev}_{E}$ is diagonal:
  \begin{equation}\label{eq:evsplit}
    \operatorname{ev}_{E}(z,v,\pi)_{i}
    =z+v\lrcorner\,\lambda_{i}+\tfrac12\,\pi\lrcorner\,\lambda_{i}^{2}
    \;\in\;H^{0,2}(A_{0})=\Ext^{2}\bigl(L_{i}[p_{i}],L_{i}[p_{i}]\bigr).
  \end{equation}
\item For every $x\in HT^{2}(A_{0})$,
  \begin{equation}\label{eq:sigmaev}
    \sigma_{E}\bigl(\operatorname{ev}_{E}(x)\bigr)=x\lrcorner\,\ch(E),
  \end{equation}
  and consequently
  $\ker\sigma_{E}\cap\image\operatorname{ev}_{E}
   =\operatorname{ev}_{E}\bigl(\{x:x\lrcorner\,\ch(E)=0\}\bigr)$.
  The weakened hypothesis of \cite[Question 11.4]{Mar25b} holds for $E$ if
  and only if $\operatorname{ev}_{E}$ kills every $x$ that preserves
  $\ch(E)$ to first order.
\item $\dim HT^{2}(A_{0})=2\binom{2n}{2}+4n^{2}$ and
  $\dim\Ext^{2}(E,E)=s\binom{2n}{2}$; for $s\ge5$ and every $n\ge2$ the
  second exceeds the first, so $\operatorname{ev}_{E}$ is not surjective and
  \cite[Lemma 11.3]{Mar25b} cannot apply to a sum of five or more line
  bundles.
\item For the object of \Cref{prop:explicitobject}, with
  $\lambda_{i}=u_{i}\theta^{+}+\bbar u_{i}\theta^{-}$ as in
  \Cref{not:splitcoords}: the space $\{x:x\lrcorner\,\ch(E)=0\}$ is
  $\Ann_{H^{0,2}}(\omega_{1})\oplus\{v:v\lrcorner\,\omega_{1}=0\}\oplus
  \{\pi:\pi\lrcorner\,\omega_{1}=0\}$, of dimensions $9$, $18$ and $9$; the
  first and the third summands are both carried by $\operatorname{ev}_{E}$
  onto $\Delta$; the second is carried onto
  \[
    \Sigma=\bigl\{\,(u_{i}a+\bbar u_{i}b)_{i=1}^{6}\;:\;
      a\in\textstyle\bigwedge^{2}V_{+}^{0,1},\ b\in\bigwedge^{2}V_{-}^{0,1}\,\bigr\},
    \qquad \dim\Sigma=n(n-1)=6,
  \]
  which contains the image of the obstruction map of
  \Cref{prop:splitobstruction}; $\Delta\cap\Sigma=0$; and
  \[
    \ker\sigma_{E}=\Delta\oplus\Sigma\subseteq\image\operatorname{ev}_{E}.
  \]
  So every direction in which $\sigma_{E}$ is blind is the value of
  $\operatorname{ev}_{E}$ on a Hochschild class that preserves $\ch(E)$ to
  first order, and the weakened hypothesis fails for $E$ by all fifteen
  dimensions. The rank of $\operatorname{ev}_{E}$ is $45$.
\end{enumerate}
\end{proposition}

\begin{proof}
(i) The evaluation map is the action of $HH^{*}(A_{0})$ on
$\Ext^{*}(E,E)$, which under HKR is contraction of a polyvector in
$H^{q}(\bigwedge^{p}T)$ with $\tfrac{1}{p!}\mathrm{At}(E)^{p}$
\cite{BF08}; for $p=0$ this is $z\cdot\id_{E}$ as in the proof of
\Cref{thm:weakfails}, for $p=1$ it is $v\lrcorner\,\mathrm{At}(E)$, the map
\eqref{eq:obdef}, and for $p=2$ it is $\tfrac12\pi\lrcorner\,\mathrm{At}(E)^{2}$.
The Atiyah class of $E$ is diagonal with entries $\lambda_{i}$, so all three
are diagonal with the entries displayed.

(ii) By \eqref{eq:splitsigma}, $\sigma_{E}(\zeta)=\sum_{i}(-1)^{p_{i}}\zeta_{i}
e^{\lambda_{i}}$ for diagonal $\zeta$. The three terms of \eqref{eq:evsplit}
are exactly the factors by which $z\wedge$, $v\lrcorner$ and $\pi\lrcorner$
multiply $e^{\lambda_{i}}$: the first trivially, the second because
$v\lrcorner$ is a derivation and $v\lrcorner\,\lambda_{i}$ has even degree,
the third by the formula for $\pi\lrcorner\,e^{\lambda}$ above. Hence
$\sigma_{E}(\operatorname{ev}_{E}(x))=\sum_{i}(-1)^{p_{i}}\,x\lrcorner\,
e^{\lambda_{i}}=x\lrcorner\,\ch(E)$. The description of
$\ker\sigma_{E}\cap\image\operatorname{ev}_{E}$ follows: $\operatorname{ev}_{E}(x)$
lies in $\ker\sigma_{E}$ if and only if $x\lrcorner\,\ch(E)=0$.

(iii) The three summands of \eqref{eq:HT2} have dimensions $\binom{2n}{2}$,
$4n^{2}$ and $\binom{2n}{2}$. The inequality
$s\binom{2n}{2}>2\binom{2n}{2}+4n^{2}$ is $(s-2)(2n-1)>4n$, and
$4n/(2n-1)\le8/3$ for $n\ge2$, so it holds for $s\ge5$.

(iv) Since $\ch(E)=12\omega_{1}$ is of type $(3,3)$, the three summands of
$x\lrcorner\,\ch(E)$ have types $(3,5)$, $(2,4)$ and $(1,3)$, so the kernel
is the direct sum of the three kernels. The first is $\Ann_{H^{0,2}}(\omega_{1})$,
of dimension $n^{2}=9$ by \Cref{lem:annihilator}. The second is
$\Hom(V_{+}^{1,0},V_{+}^{0,1})\oplus\Hom(V_{-}^{1,0},V_{-}^{0,1})$, of
dimension $2n^{2}=18$, by the proof of \Cref{prop:firstorder}, whose
argument for (i)$\Leftrightarrow$(iii) does not use the polarisation. For the
third, decompose $\bigwedge^{2}(H^{1,0})^{\vee}$ as
$\bigwedge^{2}(V_{+}^{1,0})^{\vee}\oplus(V_{+}^{1,0})^{\vee}\otimes(V_{-}^{1,0})^{\vee}
\oplus\bigwedge^{2}(V_{-}^{1,0})^{\vee}$; here $\omega_{1}$ is a combination
of $(\theta^{+})^{3}$ and $(\theta^{-})^{3}$ with both coefficients nonzero,
$(\theta^{\pm})^{3}$ spans $\bigwedge^{3}V_{\pm}^{1,0}\otimes\bigwedge^{3}V_{\pm}^{0,1}$,
the outer blocks of $\pi$ act injectively on the respective factor
$\bigwedge^{3}V_{\pm}^{1,0}$ and land in different summands of $H^{1,3}$,
and the middle block kills both; so the third kernel is the middle block,
of dimension $n^{2}=9$.

At a split point $\theta^{\pm}$ are of type $(1,1)$, so $\theta^{\pm}\in
V_{\pm}^{1,0}\otimes V_{\pm}^{0,1}$, and they are nondegenerate pairings
because $(\theta^{\pm})^{3}\ne0$. Write $\lambda_{i}^{2}=u_{i}^{2}(\theta^{+})^{2}
+\bbar u_{i}^{2}(\theta^{-})^{2}+2\Nm(u_{i})\,\theta^{+}\theta^{-}$. A middle
block $\pi$ kills $(\theta^{\pm})^{2}$, which have no $V_{\mp}^{1,0}$ slot, so
$\tfrac12\pi\lrcorner\,\lambda_{i}^{2}=\Nm(u_{i})\,\pi\lrcorner\,(\theta^{+}\theta^{-})$;
the six $u_{i}$ of \Cref{prop:explicitobject} are $\pm\zeta_{6},\pm1,
\pm\zeta_{6}^{-1}$, all of norm one, so this value is independent of $i$ and
$\operatorname{ev}_{E}(\pi)=z_{\pi}\cdot\id_{E}$ with
$z_{\pi}=\pi\lrcorner\,(\theta^{+}\theta^{-})$. The map $\pi\mapsto z_{\pi}$
from the middle block to $V_{+}^{0,1}\otimes V_{-}^{0,1}$ is the tensor
product of the two nondegenerate pairings, hence an isomorphism onto the
middle block of $H^{0,2}$, which is $\Ann_{H^{0,2}}(\omega_{1})$ by
\Cref{lem:annihilator}. So $\operatorname{ev}_{E}$ carries the third kernel
onto $\Delta$, as it carries the first by \Cref{thm:weakfails}.

For $v=(v_{+},v_{-})$ in the second kernel,
$v\lrcorner\,\lambda_{i}=u_{i}\,v_{+}\lrcorner\,\theta^{+}+\bbar u_{i}\,
v_{-}\lrcorner\,\theta^{-}$ with $a=v_{+}\lrcorner\,\theta^{+}\in
\bigwedge^{2}V_{+}^{0,1}$ and $b=v_{-}\lrcorner\,\theta^{-}\in\bigwedge^{2}V_{-}^{0,1}$;
as $v_{\pm}$ ranges over $\Hom(V_{\pm}^{1,0},V_{\pm}^{0,1})$, $a$ and $b$
range over the whole of $\bigwedge^{2}V_{\pm}^{0,1}$ because the pairings
$\theta^{\pm}$ are nondegenerate, which is the surjectivity of
$\operatorname{alt}(Vc)$ in \eqref{eq:phidelta}. So the image is $\Sigma$.
The map $(a,b)\mapsto(u_{i}a+\bbar u_{i}b)_{i}$ is injective, because
$(u_{1},\bbar u_{1})$ and $(u_{2},\bbar u_{2})$ are linearly independent in
$\CC^{2}$, $u_{1}/\bbar u_{1}=\zeta_{6}^{2}$ differing from
$u_{2}/\bbar u_{2}=1$; hence $\dim\Sigma=2\binom{n}{2}=n(n-1)=6$. The
obstruction map of \Cref{prop:splitobstruction} is the restriction to the
polarised $K$-commuting $v$, so its image lies in $\Sigma$. An element of
$\Sigma$ of the shape $(z,\dots,z)$ has $u_{i}a+\bbar u_{i}b=z=-z$ on
comparing $i$ with $3+i$, so $z=0$ and then $a=b=0$: $\Delta\cap\Sigma=0$.
By (ii), $\Delta\oplus\Sigma\subseteq\ker\sigma_{E}\cap\image\operatorname{ev}_{E}$,
and its dimension $9+6=15$ equals $\dim\ker\sigma_{E}$ by
\Cref{prop:explicitobject}, so the three spaces coincide.

For the rank, each term of \eqref{eq:evsplit} is a scalar depending on $i$
times a class independent of $i$. Write $\mathbf{u}^{k}=(u_{i}^{k})_{i}\in\CC^{6}$
and $M=V_{+}^{0,1}\otimes V_{-}^{0,1}$. The class $z$ contributes
$\mathbf{u}^{0}\otimes z$. Split $v$ into its four blocks
$v_{\varepsilon\varepsilon'}\colon V_{\varepsilon}^{1,0}\to V_{\varepsilon'}^{0,1}$;
then $v\lrcorner\,\theta^{+}$ is the sum of $v_{++}\lrcorner\,\theta^{+}\in
\bigwedge^{2}V_{+}^{0,1}$ and $v_{+-}\lrcorner\,\theta^{+}\in M$, the first
ranging over $\bigwedge^{2}V_{+}^{0,1}$ as above and the second over $M$
because $\theta^{+}$ is a nondegenerate pairing, and symmetrically for
$v\lrcorner\,\theta^{-}$ with the blocks $v_{--}$ and $v_{-+}$. As
$\bbar u_{i}=u_{i}^{-1}$, the values on the second summand of
\eqref{eq:HT2} are therefore $\mathbf{u}^{1}\otimes(a+c)+\mathbf{u}^{-1}\otimes(b+c')$
with $a$, $b$, $c$, $c'$ ranging independently over
$\bigwedge^{2}V_{+}^{0,1}$, $\bigwedge^{2}V_{-}^{0,1}$, $M$, $M$. On the
third summand, the outer block $\pi_{++}$ kills $(\theta^{-})^{2}$ and
$\theta^{+}\theta^{-}$, and $\pi_{++}\mapsto\frac12\pi_{++}\lrcorner\,(\theta^{+})^{2}$
is an isomorphism onto $\bigwedge^{2}V_{+}^{0,1}$ (in bases with
$\theta^{+}=\sum_{j}e_{j}f_{j}$ it sends $e_{j}^{\vee}\wedge e_{k}^{\vee}$ to
$\pm f_{j}f_{k}$); likewise for
$\pi_{--}$, and the middle block contributes $\mathbf{u}^{0}\otimes z_{\pi}$.
So
\[
  \image\operatorname{ev}_{E}
  =\mathbf{u}^{0}\otimes H^{0,2}
  +\mathbf{u}^{\pm1}\otimes\bigl(\textstyle\bigwedge^{2}V_{\pm}^{0,1}\oplus M\bigr)
  +\mathbf{u}^{\pm2}\otimes\textstyle\bigwedge^{2}V_{\pm}^{0,1}.
\]
The $u_{i}$ are the six sixth roots of unity, so the vectors
$\mathbf{u}^{k}$ for $k=0,\pm1,\pm2$ are distinct characters of the cyclic
group they form and are linearly independent. The sum is direct, and the
rank is $15+2\cdot12+2\cdot3=45$.
\end{proof}

\Cref{fig:evaluation} draws the three kernels and their images in
$\Ext^{2}(E,E)$.

\begin{figure}[tbp]
\centering
  \includegraphics[width=\textwidth]{fig_evaluation}
\caption{\Cref{prop:evaluation} for the explicit object at $n=3$. Left: the
three summands of $HT^{2}(A_{0})$, drawn with heights proportional to their
dimensions, and inside each the part that preserves $\ch(E)$ to first order:
the annihilator of $\omega_{1}$, the deformations killing $\omega_{1}$ with
the polarised $K$-commuting tangent space $T$ in front, and the middle block
of bivectors. Right: $\Ext^{2}(E,E)$ as six stacked copies of $H^{0,2}$, one
per summand of $E$, and inside it the kernel of $\sigma_{E}$: the column
$\Delta$, one and the same class in every layer, and the column $\Sigma$,
whose entries depend on the summand through $u_{i}$. The three arrows are the
evaluation map on the three kernels; the two teal ones both land on
$\Delta$, the ochre one on $\Sigma$, and together they cover the kernel.
Below, the commutative square that makes this a statement about the weakened
criterion of \cite[Question 11.4]{Mar25b}.}
\label{fig:evaluation}
\end{figure}

\begin{remark}[Question 11.4 for split objects]
\label{rem:question114}
For a split object the two criteria stand or fall together, and for the
explicit object both fall completely. \Cref{thm:weakfails} makes the first
$n^{2}$ dimensions of the failure uniform over every concentrated object in
every dimension; \Cref{prop:evaluation}(iv) accounts for the remaining
$n(n-1)$ at $n=3$ by the deformations of the two halves $V_{\pm}$ of the
Hodge structure separately, and the equality $\ker\sigma_{E}=\Delta\oplus\Sigma$
says that the two mechanisms exhaust the kernel: nothing in it is
accidental. None of this touches the case for which the question was asked.
A sheaf of positive rank has $z\cdot\ch(E)\ne0$ for every $z\ne0$, so its
$\Delta$ is zero, and whether $\operatorname{ev}_{E}$ kills the deformations
that preserve $\kappa(E)$ is a statement about the sheaf and not about its
Chern character. Markman's candidate in \cite[Section 12]{Mar25b}, a rank
one complex on a principally polarised abelian fourfold with Chern character
on the secant $\langle e^{\pm\sqrt{-d}\Theta}\rangle$, is exactly such a
sheaf, and he records that the verification remains to be done.
\Cref{thm:factor} reduces that verification to one identity on the fourfold,
and \Cref{thm:candidatefails} shows that the identity fails: the candidate
does not satisfy the weakened criterion either.
\end{remark}

\subsubsection*{Solutions of nonzero rank}

We return to the Diophantine problem of \Cref{setup:pte}. By
\Cref{cor:nplusone} a solution with $\ch_{0}(E)\ne0$ needs at least $n+1$
distinct norms, and multiplicity one bounds each level sum by the number of
elements of $\cO_{K}$ of that norm.

\begin{proposition}[The level sums]\label{thm:pterange}
Keep \Cref{setup:pte} and suppose $E$ satisfies \eqref{eq:pteconditions} with
$\ch_{0}(E)\ne0$, so that the $u_{i}$ take $\nu\ge n+1$ distinct norms by
\Cref{cor:nplusone}. Write $N_{1}<\dots<N_{\nu}$ for those norms, $M_{j}$ for
the sum of the $m_{i}$ at level $N_{j}$, and
$r_{K}(N)=\#\{u\in\cO_{K}:\Nm_{K/\QQ}(u)=N\}$. If $\nu=n+1$ then, with
$P=\prod_{k}N_{k}$, $S=\ch_{0}(E)$ and
$D_{j}=N_{j}\prod_{k\ne j}(N_{j}-N_{k})$,
\begin{equation}\label{eq:levelsums}
  M_{j}=(-1)^{n}\,\frac{S\,P}{D_{j}},
\end{equation}
the least $S>0$ making every $M_{j}$ an integer is
$S_{\min}=\operatorname{lcm}_{j}\bigl(|D_{j}|/\gcd(|D_{j}|,P)\bigr)$, and
multiplicity one, \Cref{lem:splitsemireg}(i), forces
$|M_{j}|\le r_{K}(N_{j})$ for every $j$.
\end{proposition}

\begin{proof}
The conditions $T_{r,r}=0$ for $r=1,\dots,n$ say that the linear
functional $f\mapsto\sum_{j}M_{j}f(N_{j})$ kills $x,x^{2},\dots,x^{n}$, hence
sends every polynomial $f$ of degree at most $n$ to $S\,f(0)$. Applying it to
$f(x)=\prod_{k\ne j}(x-N_{k})$, which has degree $n$, gives
\[
  M_{j}\prod_{k\ne j}(N_{j}-N_{k})=S\,(-1)^{n}\prod_{k\ne j}N_{k},
\]
and multiplying by $N_{j}$ turns this into \eqref{eq:levelsums}. The value of
$S_{\min}$ is the definition of a least common multiple, and the bound on
$|M_{j}|$ holds because each $u$ of norm $N_{j}$ contributes at most one term
$m_{i}=\pm1$ to $M_{j}$.
\end{proof}

So the level sums are determined by the norms up to the scale $S$, and they
compete with the supply $r_{K}(N_{j})$ of elements at each level. That
competition is arithmetic. A second mechanism, in the cohomology of $A$,
closes the split route for every $n$ and every set of norms at once.

\begin{corollary}[Sums of line bundles fail both criteria]\label{cor:noline}
Let $A$ be of $(K,-1,n)$-Weil type with $n\ge2$, and let
$E=\bigoplus_{i}L_{i}[p_{i}]^{\oplus m_{i}}$ be a sum of shifted line bundles,
twisted by a common class or not, whose (twisted) Chern character remains of
Hodge type to first order along $\cD_{n,n}$ and has a nonzero Weil component.
Then $\sigma_{E}$ is not injective on the image of $\operatorname{ev}_{E}$:
$E$ is neither semiregular nor semiregular in the weakened sense of
\cite[Question 11.4]{Mar25b}, and it is obstructed to first order along some
direction of the Weil family. In particular every object of \Cref{setup:pte}
satisfying \eqref{eq:pteconditions}, for every $n\ge2$, every field $K$,
every rank and every set of norms, fails both criteria.
\end{corollary}

\begin{proof}
Write $\lambda_{i}$ for the twisted first Chern classes of the summands. If
every $\lambda_{i}$ were a real multiple of $\eta$, it would be a rational
multiple, both classes being rational, and the twisted Chern character
$\sum_{i}m_{i}e^{\lambda_{i}}$ would be a polynomial in $\eta$; its component
in degree $2n$ would then lie in $\QQ\eta^{n}$, which meets the Weil line in
zero by \Cref{prop:grading}, against the hypothesis. So some $\lambda_{i}$ is
not a real multiple of $\eta$, and \Cref{thm:linebundle} applies to the
summand $L_{i}[p_{i}]$. For an object of \Cref{setup:pte} the classes are
$\lambda_{i}=\theta(u_{i})\in R_{0}$, and $R_{0}\cap\QQ\eta=0$; no twist is
needed because $c_{1}(E)=\theta(T_{1,0})=0$ by \eqref{eq:pteconditions}, as
$(1,0)$ is not an excluded pair for $n\ge2$; and \eqref{eq:pteconditions}
gives a flat Chern character with a nonzero Weil component by
\Cref{prop:chintermsofu}.
\end{proof}

\begin{remark}[What the corollary asks of an object]\label{rem:nolinemeaning}
By \Cref{cor:noline} there is no sum of line bundles, at any member of any
Weil family in any dimension, that could run the semiregularity argument in
either form. The reason is \Cref{prop:rigidity}(i): the
only divisor class that stays of type $(1,1)$ across the whole family is the
polarisation, so every line bundle summand off that line is obstructed in
some direction of the family, and the semiregularity map, whose composite
with the obstruction is contraction with the flat class $\ch(E)$, cannot see
it. The object the argument needs is one with no rank one
summand off the polarisation line, of positive rank, indecomposable enough
that its Atiyah class contracted with every direction of $T$ vanishes.
Markman's reflexive sheaves at $n=3$ are of that kind, and nothing here
says that their analogues at $n\ge4$ do not exist.
\end{remark}

\subsection{Convolutions of line bundles}\label{ssec:convolutions}

After \Cref{cor:noline} the next objects built from line bundles are their
convolutions, the iterated cones of shifted line bundles. Their Chern
character is the same signed sum of exponentials as for a direct sum, but
their $\Ext$ groups are cut down by the maps between the pieces, so the
argument of \Cref{thm:linebundle} no longer applies as it stands. This
subsection shows that fewer than $2n$ pieces with nondegenerate differences
never serve (\Cref{thm:fewpieces}); writes down, at a CM member for
$K=\QQ(i)$, a signed sum of $2\cdot4^{n-1}$ line bundles that is exactly a
real Weil class (\Cref{prop:weilpieces}); and determines, for the
convolutions of those pieces with three multiples of the polarisation at
$n=3$, which products in the twisted differential can act on the $525$
diagonal classes of degree two, of which at least $468$ must disappear
before the polarised criterion can hold (\Cref{thm:esixdiagonal}). Nothing
enters the diagonal, cup products remove at most $276$ of these classes, and
everything else would have to be removed by fourfold products along a path
through the three multiples and one of the Weil pieces: one family of them
for some placements of the multiples, and two families that exclude each
other for the others (\Cref{prop:mixedplacement}); one order of the pieces
gives a convolution that solves the Maurer--Cartan equation
(\Cref{prop:explicitconvolution}). Off the diagonal, the classes between two
pieces that differ in every coordinate are reached by no product
(\Cref{lem:allthree}), and no
convolution of these pieces meets the polarised criterion
(\Cref{thm:noconvolution}). At $n=4$ the same classes, now between pieces
that differ in three coordinates, are reached by no product in any
convolution, and they exclude every convolution whose $128$ pieces
$L_{\zeta}$ occupy two shifts (\Cref{thm:efourtwolevels}) or three
consecutive ones (\Cref{thm:efourthreelevels}); with the groups between
pieces whose shifts differ by two or three (\Cref{lem:efourspread}) they
exclude every convolution whose pieces have their shifts within five
consecutive values (\Cref{cor:efourfive}); the diagonal classes at a lonely
shift (\Cref{lem:efourlonely}) and the interleaved shifts
(\Cref{prop:efourpairs}) extend this to six (\Cref{cor:efoursix}).

\begin{setupenv}[Convolutions]\label{setup:convolution}
Let $A$ be an abelian variety of dimension $g$ with a flat K\"ahler metric,
and let $L_{1},\dots,L_{m}$ be line bundles on $A$ with hermitian metrics of
constant curvature, hermitian forms $H_{1},\dots,H_{m}$ and pairwise distinct
first Chern classes. For integers $d_{1},\dots,d_{m}$ put
$P_{a}=L_{a}[-d_{a}]$. A \emph{convolution} of $P_{1},\dots,P_{m}$ is an
object $E$ given by a one-sided twisted
complex\footnote{Every object obtained from $P_{1},\dots,P_{m}$ by iterated
cones in this order is of this form, because the Dolbeault complexes form a
differential graded enhancement of the full subcategory of line bundles. All
the objects of this subsection are given in this way.}
\[
  \Bigl(\,\bigoplus_{a}\cA^{0,*}(L_{a})[-d_{a}],\ \bar\partial+\delta\Bigr),
  \qquad
  \delta=(\delta_{ab})_{a<b},\quad
  \delta_{ab}\in\cA^{0,1+d_{a}-d_{b}}\bigl(L_{a}^{-1}\otimes L_{b}\bigr),
\]
with $\bar\partial\delta+\delta\circ\delta=0$, so that
$\ch(E)=\sum_{a}(-1)^{d_{a}}e^{c_{1}(L_{a})}$. Write
$\cA(a,b)=\cA^{0,*}(L_{a}^{-1}\otimes L_{b})$ and $\cH(a,b)$ for its harmonic
forms, so that $\cH^{q}(a,b)\cong H^{q}(A,L_{a}^{-1}\otimes L_{b})$, and write
$\pi$, $G$ and $h=\bar\partial^{*}G$ for the harmonic projection, the Green
operator and the homotopy, with $1-\pi=\bar\partial h+h\bar\partial$ and
$\pi h=h\pi=h^{2}=0$. A form of degree $q$ in $\cA(a,b)$ is a morphism
$P_{a}\to P_{b}$ of degree $w=q-d_{a}+d_{b}$; we call $q$ its $\Ext$ degree
and, when $w=1$, call $q-1$ its \emph{excess}. When $H_{b}-H_{a}$ is
nondegenerate, $\cH(a,b)$ is concentrated in the single degree
$n(H_{b}-H_{a})$, the number of negative eigenvalues \cite[Chapter~3]{BL04}.
\end{setupenv}

\begin{lemma}[The harmonic model]\label{lem:harmonicmodel}
In \Cref{setup:convolution}:
\begin{enumerate}[label=\textup{(\roman*)},leftmargin=2.6em,itemsep=2pt]
\item $E$ is isomorphic to a convolution of the same pieces with
  $h\delta=0$, and then $\delta=\delta^{H}-h(\delta\circ\delta)$ with
  $\delta^{H}=\pi\delta$ harmonic. Substituting repeatedly, every
  $\delta_{ab}$ is a finite sum of planar binary trees whose leaves are
  components $\delta^{H}_{cd}$ with $c<d$, whose internal vertices compose,
  and whose internal edges carry $-h$.
\item $\Ext^{*}(E,E)$ is the cohomology of
  $\cH_{E}=\bigoplus_{a,b}\cH(a,b)$, graded by $w$, with the differential
  $d_{E}=\sum_{k\ge0}\pi\,Q(-hQ)^{k}$, where $Q=[\delta,-]$.
\item For $x\in\cH(a,b)$, the component of $d_{E}(x)$ in $\cH(c,e)$ is a sum
  of terms $\pm\pi\,T$, one for each planar binary tree $T$ with $k\ge2$
  leaves, internal vertices composing, $-h$ on the internal edges other
  than the root, whose leaves read in order are the components along a path
  $c=a_{0}\to a_{1}\to\dots\to a_{k}=e$ in which one step is
  $x\colon a\to b$ and the others are components $\delta^{H}_{a_{j}a_{j+1}}$
  with $a_{j}<a_{j+1}$. The $\Ext$ degree of such a term is the sum of the
  $\Ext$ degrees of its leaves minus $k-2$.
\end{enumerate}
\end{lemma}

The terms with $k$ leaves in (iii) are the Massey products of length $k$ of
the pieces;\footnote{In the language of $A_{\infty}$-categories they are the
higher products $m_{k}$ of the minimal model of the line bundles, evaluated
on the Maurer--Cartan element $\delta^{H}$ of the twisted complex; for
$k=2$ they are cup products.} (iii) is all that is used of them below.

\begin{proof}
(i) Suppose $h\delta_{cd}=0$ for $d-c<r$. Conjugating by $1+g$, with $g$ of
degree $0$ supported at distance $r$, replaces $\delta_{ab}$ for $b-a=r$ by
$X_{ab}+\bar\partial g_{ab}$, where $X_{ab}$ depends only on $\delta$, and
leaves the components at smaller distance unchanged. With
$g_{ab}=-hX_{ab}$ the new component is $X_{ab}-\bar\partial hX_{ab}$, and
$h\bar\partial h=h(1-\pi-h\bar\partial)=h$ gives $h(X_{ab}-\bar\partial
hX_{ab})=0$. Induction on $r$ gives $h\delta=0$, and then
$\delta=\pi\delta+\bar\partial h\delta+h\bar\partial\delta
=\delta^{H}-h(\delta\circ\delta)$ by the Maurer--Cartan equation. The
component $(\delta\circ\delta)_{ab}$ involves only components at distance
less than $b-a$, so the substitution ends.
(ii) The endomorphism complex of $E$ is $\bigoplus_{a,b}\cA(a,b)$ with the
differential $\bar\partial+Q$. Blockwise, the inclusion of harmonic forms,
$\pi$ and $-h$ form a strong deformation retraction of
$(\bigoplus\cA(a,b),\bar\partial)$ onto $(\cH_{E},0)$, and $Q$ raises $b-a$,
so $hQ$ is nilpotent. The basic perturbation lemma gives a homotopy
equivalence of $(\bigoplus\cA(a,b),\bar\partial+Q)$ with $(\cH_{E},d_{E})$,
with the chain maps $\iota_{E}=\sum_{k}(-hQ)^{k}$ on $\cH_{E}$ and
$\pi_{E}=\pi\sum_{k}(-Qh)^{k}$. (iii) Each $Q$ in $d_{E}$ composes with one
component of $\delta$ on the left or on the right; writing that component as
in (i) gives the trees, the edges from (i) and those from $d_{E}$ being
exactly the internal edges other than the root. Composition adds degrees,
$h$ lowers them by one, and a binary tree with $k$ leaves has $k-2$ internal
edges.
\end{proof}

\begin{lemma}[The index of a sum]\label{lem:indexsum}
For hermitian forms $X,Y$ on $\CC^{g}$, $n(X+Y)\le n(X)+n(Y)$. If
$H_{1}+\dots+H_{k}=0$, then $\sum_{j}n(H_{j})\ge g-\dim\ker H_{i}$ for every
$i$; in particular $\sum_{j}n(H_{j})\ge g$ when one $H_{i}$ is
nondegenerate.
\end{lemma}

\begin{proof}
A subspace on which $X+Y$ is negative definite meets the subspaces on which
$X\ge0$ and $Y\ge0$, of codimensions $n(X)$ and $n(Y)$, only in zero, so its
dimension is at most $n(X)+n(Y)$. For the second statement,
$-H_{i}=\sum_{j\ne i}H_{j}$ gives
$g-n(H_{i})-\dim\ker H_{i}=n(-H_{i})\le\sum_{j\ne i}n(H_{j})$.
\end{proof}

\begin{theorem}[Fewer than $2n$ pieces]\label{thm:fewpieces}
Let $A$ be of $(K,-1,n)$-Weil type with $n\ge2$, and let $E$ be a
convolution of $m<2n$ shifted line bundles whose differences $H_{b}-H_{a}$
are all nondegenerate. Suppose that $\ch(E)$ remains of Hodge type to first
order along $\cD_{n,n}$ and has a nonzero Weil component, as every polarised
character $N\omega_{1}+p(\theta)$ with $N\ne0$ does. Then there is
$v\in T_{[A]}\cD_{n,n}$ with
\[
  \operatorname{ev}_{E}(v)\ne0
  \qquad\text{and}\qquad
  \sigma_{E}\bigl(\operatorname{ev}_{E}(v)\bigr)=0 ,
\]
so $\sigma_{E}$ is not injective, and for a polarised character
$\dim\Ext^{2}(E,E)>r(\ch E)$ by \Cref{thm:p2primenumber}(iv).
\end{theorem}

\begin{proof}
First, no element of degree one of $\cH_{E}$ has a differential with a
nonzero component in the diagonal classes $\cH^{2}(c,c)=H^{2}(A,\cO_{A})$.
Such a component is a sum of terms of \Cref{lem:harmonicmodel}(iii) whose
$k$ leaves have degree one and run around a closed path $c\to\dots\to c$; the
pieces before the step $x$ increase from $c$ and those after it increase to
$c$, so the path passes through $k$ distinct pieces and $k\le m$. Its
leaves have $\Ext$ degrees $q_{j}=n(H_{a_{j+1}}-H_{a_{j}})$, and the output
has $\Ext$ degree $\sum_{j}q_{j}-(k-2)=2$, so $\sum_{j}q_{j}=k<2n$. The $k$
differences are nondegenerate and add up to zero, so \Cref{lem:indexsum}
gives $\sum_{j}q_{j}\ge g=2n$, a contradiction. Hence taking diagonal
components gives a well defined map
$\pi_{D}\colon\Ext^{2}(E,E)\to\bigoplus_{a}H^{2}(A,\cO_{A})$.

Next, $\operatorname{ev}_{E}(v)=v\lrcorner\mathrm{At}(E)$, as in the proof
of \Cref{thm:linebundle}. In the twisted complex the Atiyah class is
represented by the commutator of the Chern $(1,0)$-connection
$\bigoplus_{a}\nabla_{a}$ with $\bar\partial+\delta$; its diagonal entries
are the curvatures $F_{a}$, constant forms representing $c_{1}(L_{a})$, and
its other entries $\nabla\delta_{ab}$ lie above the diagonal. So
$\operatorname{ev}_{E}(v)$ is represented by an upper triangular $z$ with
$z_{aa}=v\lrcorner F_{a}$, and in $\cH_{E}$ by $\pi_{E}(z)$. The operators
$Q$ and $h$ preserve upper triangular forms and $Q$ makes them strictly upper
triangular, so the diagonal of $\pi_{E}(z)$ is that of $\pi(z)$, and
$\pi_{D}(\operatorname{ev}_{E}(v))=(v\lrcorner c_{1}(L_{a}))_{a}$.

If every $c_{1}(L_{a})$ were a real multiple of the polarisation, it would
be a rational one and $\ch(E)$ a polynomial in it, whose component in degree
$2n$ meets the Weil line in zero by \Cref{prop:grading}. So some
$c_{1}(L_{a})$ is not, and by \Cref{prop:rigidity}(i) some
$v\in T_{[A]}\cD_{n,n}$ has $v\lrcorner c_{1}(L_{a})\ne0$; then
$\pi_{D}(\operatorname{ev}_{E}(v))\ne0$, while
$\sigma_{E}(\operatorname{ev}_{E}(v))=v\lrcorner\ch(E)=0$.
\end{proof}

\Cref{thm:linebundle} is the case $\delta=0$, with any number of pieces. A
convolution reaches its diagonal classes once it has $2n$ pieces, when a
cycle through all of them can have total index $2n$, the equality case of
\Cref{lem:indexsum}. A convolution meeting the polarised criterion must
therefore have at least $2n$ pieces, or pieces with degenerate differences.
The next construction has both.

\begin{proposition}[A pure Weil class from line bundles]
\label{prop:weilpieces}
Let $E_{0}=\CC/\ZZ[i]$, $A=E_{0}^{2n}$ with coordinates $z_{1},\dots,z_{2n}$
and $n\ge2$, let $K=\QQ(i)$ act by $i$ on $z_{1},\dots,z_{n}$ and by $-i$
on $z_{n+1},\dots,z_{2n}$, and let $\theta$ be the product principal
polarisation; then $A$ is of $(\QQ(i),-1,n)$-Weil type, a CM member as in
\Cref{rem:cmquadratic}. Write $B_{j}$ for the surface with coordinates
$(z_{j},z_{n+j})$, so that $A=B_{1}\times\dots\times B_{n}$. For $p\in\ZZ$
and $\zeta\in\mu_{4}^{n}$ let $L_{\zeta}$ be a line bundle whose hermitian
form is $\bigl(\begin{smallmatrix}p&\zeta_{j}\\ \bbar\zeta_{j}&p
\end{smallmatrix}\bigr)$ on $B_{j}$ and zero between different surfaces, and
put $\alpha=\prod_{j}\tfrac{i}{2}\,dz_{j}\wedge d\bar z_{n+j}$.
\begin{enumerate}[label=\textup{(\roman*)},leftmargin=2.6em,itemsep=2pt]
\item $\alpha$ and $\bbar\alpha$ span the Weil classes of $A$, and
  \[
    \sum_{\prod\zeta=1}e^{c_{1}(L_{\zeta})}
    -\sum_{\prod\zeta=-1}e^{c_{1}(L_{\zeta})}
    =2\cdot4^{n-1}\,(\alpha+\bbar\alpha).
  \]
\item For distinct integers $t_{1},t_{2},t_{3}$ and signs $\sigma_{i}$, the
  character
  $\gamma=2\cdot4^{n-1}(\alpha+\bbar\alpha)+\sum_{i}\sigma_{i}e^{t_{i}\theta}$
  has Hankel rank $3$, so $r(\gamma)=7n^{2}-2n$. A convolution with this
  character of the $2\cdot4^{n-1}$ bundles $L_{\zeta}$, with
  $(-1)^{d}=\prod\zeta$, and of three line bundles $M_{i}$ with
  $c_{1}(M_{i})=t_{i}\theta$ has $2\cdot4^{n-1}+3$ pieces and
  $(2\cdot4^{n-1}+3)\,n(2n-1)$ diagonal classes of degree two: $525$
  against $r=57$ at $n=3$, and $3668$ against $104$ at
  $n=4$.\footnote{At $n=2$ the same construction has $8+3=11$ pieces and
  $66$ diagonal classes against $r=24$.}
\end{enumerate}
\end{proposition}

\begin{proof}
The hermitian forms have imaginary parts integral on $\ZZ[i]^{2n}$, so the
line bundles exist. On $B_{j}$ write $c_{1}$ of the restriction as
$p(a_{j}+b_{j})+\zeta_{j}c_{j}+\bbar\zeta_{j}c'_{j}$, where
$a_{j}=\tfrac i2dz_{j}\wedge d\bar z_{j}$,
$b_{j}=\tfrac i2dz_{n+j}\wedge d\bar z_{n+j}$,
$c_{j}=\tfrac i2dz_{j}\wedge d\bar z_{n+j}$ and
$c'_{j}=\tfrac i2dz_{n+j}\wedge d\bar z_{j}$. Since $c_{j}^{2}=c_{j}'^{2}=0$,
\[
  e^{c_{1}(L_{\zeta})}=\prod_{j}f_{j}(\zeta_{j}),\qquad
  f_{j}(z)=e^{p(a_{j}+b_{j})}\bigl(1+zc_{j}+\bbar zc'_{j}+c_{j}c'_{j}\bigr).
\]
For
$u\in\mu_{4}$, $\tfrac12(u+\bbar u)$ is $u$ when $u=\pm1$ and $0$ when
$u=\pm i$, so the left side of (i) is
$\tfrac12\bigl(\prod_{j}\sum_{z}zf_{j}(z)+\prod_{j}\sum_{z}\bbar zf_{j}(z)\bigr)$.
Now $\sum_{z\in\mu_{4}}z=\sum z^{2}=0$ and $\sum|z|^{2}=4$, so
$\sum_{z}zf_{j}(z)=4e^{p(a_{j}+b_{j})}c'_{j}=4c'_{j}$, because $a_{j}c'_{j}$
and $b_{j}c'_{j}$ repeat a differential, and likewise
$\sum_{z}\bbar zf_{j}(z)=4c_{j}$. This gives
$2\cdot4^{n-1}(\prod c'_{j}+\prod c_{j})=2\cdot4^{n-1}(\bbar\alpha+\alpha)$.
The differentials $dz_{j}$, $j\le n$, and $d\bar z_{n+j}$ span the
$i$-eigenspace of $K$ on $H^{1}(A,\CC)$, and $\alpha$ spans its top exterior
power, which is a Weil class, and $\bbar\alpha$ the conjugate one. (ii) The
moments $\mu_{m}=\sum_{i}\sigma_{i}t_{i}^{m}$ make the Hankel matrix of
\Cref{setup:p2prime} a product of Vandermonde matrices of rank $3$, and
\Cref{thm:p2primenumber}(ii) gives $r$. The counts are immediate.
\end{proof}

The pieces $L_{\zeta}$ and $L_{\zeta'}$ have a degenerate difference as soon
as $\zeta$ and $\zeta'$ share a coordinate, so \Cref{thm:fewpieces} does not
apply to them, and there are more than $2n$ of them. The rest of this
subsection computes what the Massey products can do to the diagonal classes
of such a convolution. The pieces are external products over the surfaces
$B_{j}$, and the first tool is a bound on how far a product can lower the
degree on each surface.

\begin{lemma}[Degree drop on a product]\label{lem:degreedrop}
Let $A=B_{1}\times\dots\times B_{n}$ with a product metric, and let every
$L_{a}$ be an external product with a product metric. Decompose the leaves
of a term of \Cref{lem:harmonicmodel}(iii) with $k$ leaves into tensor
products of harmonic forms on the factors. Then the output has degree
$s_{j}-d_{j}$ on $B_{j}$, where $s_{j}$ is the sum of the degrees of the
leaves on $B_{j}$ and $d_{j}\ge0$, $\sum_{j}d_{j}=k-2$, and
\[
  d_{j}\le\max(0,\,u_{j}-2),
\]
where $u_{j}$ is the number of leaves whose component on $B_{j}$ is not a
multiple of the unit section of $\cO_{B_{j}}$ (such a unit occurs only where
source and target agree on $B_{j}$).
\end{lemma}

\begin{proof}
The Laplacian is $\Delta=\sum_{j}\Delta_{j}$, so $h=\sum_{j}h_{j}$ with
$h_{j}=\bar\partial_{j}^{*}G=G\bar\partial_{j}^{*}$, which lowers the degree
on $B_{j}$ by one and keeps the others; expand every internal edge
accordingly, and let $d_{j}$ be the number of edges carrying $h_{j}$. Let
$V_{j}$, $U_{j}$ and $W_{j}$ be the forms whose component on $B_{j}$ is
$\Delta_{j}$-harmonic, a multiple of the unit, or in the image of
$\bar\partial_{j}^{*}$. The operators $G$ and $h_{i}$, $i\ne j$, preserve all
three, $h_{j}$ kills $V_{j}$ and $W_{j}$, $\pi$ kills $W_{j}$, and
multiplication by an element of $U_{j}$ preserves $V_{j}$, $U_{j}$ and
$W_{j}$. For an edge $e$ carrying $h_{j}$ let $N(e)$ be the set of leaves
below $e$ whose component on $B_{j}$ is not a unit. If $|N(e)|\le1$ the
value below $e$ lies in $V_{j}$ and the term vanishes. If $e$ lies below the
next such edge $e'$ with $N(e)=N(e')$, everything composed in between lies
in $U_{j}$, the value at $e'$ lies in $W_{j}$, and the term vanishes; if
$N(e)$ is the set of all $u_{j}$ such leaves for the highest such edge, the
value at the root lies in $W_{j}$ and $\pi$ kills it. So in a nonzero term
the sets $N(e)$ are distinct, of size at least two, different from the whole
set, and pairwise nested or disjoint. Together with the whole set they form
a family of that kind, which has at most $u_{j}-1$ members, since its members
are the internal vertices of a forest with $u_{j}$ leaves in which every
internal vertex has at least two children. Hence $d_{j}\le u_{j}-2$.
\end{proof}

\begin{proposition}[Nothing enters the diagonal]\label{prop:nothingenters}
Keep \Cref{prop:weilpieces} with $n\ge3$, and let $t_{1},t_{2},t_{3}$ be
distinct with $|t_{i}-p|\ge2$. For every convolution $E$ of the pieces
$L_{\zeta}$ and $M_{i}$, in any order, with any shifts of the given parities
and with the pieces twisted by any line bundles of degree zero, no element of
degree one of $\cH_{E}$ has a differential with a nonzero diagonal
component. Consequently
\[
  \dim\Ext^{2}(E,E)\;\ge\;\dim\ker\bigl(d_{E}|_{D^{2}}\bigr),
  \qquad
  D^{2}=\bigoplus_{a}H^{2}(A,\cO_{A}).
\]
\end{proposition}

\begin{proof}
As in the proof of \Cref{thm:fewpieces}, a term that reaches the diagonal
runs around a closed path of distinct pieces whose $k$ steps have degree one
and total excess $0$. The step from $L_{\zeta}$ to $L_{\zeta'}$ has $\Ext$
degree $1$ on each surface where $\zeta_{j}\ne\zeta'_{j}$, the difference
there having eigenvalues $\pm|\zeta'_{j}-\zeta_{j}|$, and at least $0$
elsewhere. If $\prod\zeta=\prod\zeta'$ the two differ on at least two
surfaces; otherwise the step has even $\Ext$ degree, the shifts having
opposite parities. Either way its excess is at least $1$. The difference
between $L_{\zeta}$ and $M_{i}$ has eigenvalues $t_{i}-p\pm1$ on every
surface, so it is definite: the step has $\Ext$ degree $0$ and excess $-1$
if it goes up the value $\phi$, where $\phi(L_{\zeta})=p$ and
$\phi(M_{i})=t_{i}$, and $\Ext$ degree $2n$ and excess $2n-1$ if it goes
down; the same holds between the $M_{i}$. A maximal run of the cycle through
the $M_{i}$ starts and ends at the value $p$, or is the whole cycle, so it
has at least one step down and at most three steps up; its excess is at least
$2n-1-3\ge2$. So every cycle has excess at least $1$. Twisting by line
bundles of degree zero leaves the forms unchanged and can only make some
$\cH(a,b)$ vanish. For the inequality, the diagonal components of cocycles
of degree two vanish on coboundaries, so $\dim\Ext^{2}(E,E)$ is at least the
dimension of their image, which contains the kernel of $d_{E}$ on $D^{2}$.
\end{proof}

At $n=2$ the bound on a run is $0$ and a cycle through the three multiples
can reach the diagonal; this is why the analysis below starts at $n=3$.

\begin{figure}[tbp]
\centering
  \includegraphics[width=\textwidth]{fig_convolutions}
\caption{The pieces of \Cref{prop:weilpieces} at their values $\phi$, with
$\phi(L_{\zeta})=p$ and $\phi(M_{i})=t_{i}$. (a) A step of degree one
between pieces with a definite difference has $\Ext$ degree $0$ going up
and $2n$ going down, so excess $-1$ or $2n-1$; a step between two
$L_{\zeta}$ has excess at least $1$. The cycle drawn runs up through all
three $M_{i}$ and down once, the worst case, and still has excess
$2n-3>0$, so for $n\ge3$ no term reaches the diagonal
(\Cref{prop:nothingenters}). (b) The one path whose products act on the
diagonal at $n=3$ when the $t_{i}$ lie above $p$
(\Cref{thm:esixdiagonal}(i)), drawn at the same heights
and in the order of the convolution, with the $\Ext$ degrees of its steps
and the shifts they force; the dashed arrow is the value of the product with
a diagonal class $y$.}
\label{fig:convolutions}
\end{figure}

\begin{proposition}[Cup products on the diagonal]\label{prop:cupkernel}
Keep \Cref{prop:nothingenters} with $n=3$, so that $D^{2}$ has dimension
$35\cdot15=525$. The terms of $d_{E}$ on $D^{2}$ with two leaves, the cup
products, vanish on the $45$ diagonal classes of the $M_{i}$, and on the
$480$ diagonal classes of the $L_{\zeta}$ they have a kernel of dimension at
least $204$. So their kernel on $D^{2}$ has dimension at least $249$.
\end{proposition}

\begin{proof}
For $y=(y_{a})\in D^{2}$ the cup products in $\cH(a,b)$ are
$(y_{a}-y_{b})\wedge\delta^{H}_{ab}$, since $y$ has even degree. On a surface
where the pieces differ, $\delta^{H}_{ab}$ has degree $1$, the only degree of
the cohomology there, and the product with a component of $y$ of positive
degree lands in a degree where the cohomology vanishes. The pieces $M_{i}$
differ from every other piece on every surface, which gives the first
statement. Write $y_{a}$ by its type $\tau\in\{0,1,2\}^{3}$, $|\tau|=2$, the
degrees on the three surfaces: three types $(2,0,0)$ of dimension one and
three types $(1,1,0)$ of dimension four. The degree count on each surface
and the parity of $\Ext$ degrees show that $L_{\zeta}$ and $L_{\zeta'}$ can
carry a nonzero cup product on the type $\tau$ only if either $\zeta'$
differs from $\zeta$ by a sign in one coordinate outside the support of
$\tau$, or $\zeta'$ is $\zeta$ multiplied by $(i,i)$ or $(-i,-i)$ in two
coordinates and $\tau$ is $2$ on the third. Let $\Gamma_{\tau}$ be the graph
on the $32$ pieces with these edges. If each component $y^{\tau}$ is
constant on the connected components of $\Gamma_{\tau}$, every cup product
vanishes. For $\tau=(2,0,0)$ the coordinate $\zeta_{1}$ is constant along
the edges, and the eight pieces with a given $\zeta_{1}$ are connected, so
$\Gamma_{\tau}$ has $4$ components; for $\tau=(1,1,0)$ the edges change the
sign of $\zeta_{3}$ only, and $\Gamma_{\tau}$ has $16$ components. The
kernel therefore contains a space of dimension $3\cdot4\cdot1+3\cdot16\cdot4
=12+192=204$.
\end{proof}

\begin{theorem}[The diagonal on $E_{0}^{6}$]\label{thm:esixdiagonal}
Let $n=3$ and let $E$ be a convolution of the $32$ bundles $L_{\zeta}$ and
of $M_{1},M_{2},M_{3}$ as in \Cref{prop:weilpieces}, with
$p+2\le t_{1}<t_{2}<t_{3}$ and signs $\sigma_{i}$, so that
$\ch(E)=32(\alpha+\bbar\alpha)+\sum_{i}\sigma_{i}e^{t_{i}\theta}$ and the
polarised criterion asks for $\dim\Ext^{2}(E,E)=57$.
\begin{enumerate}[label=\textup{(\roman*)},leftmargin=2.6em,itemsep=2pt]
\item The Massey products of length at least three vanish on $D^{2}$, except
  those of length four along a path
  $M_{2}\to M_{3}\to L_{\zeta}\to M_{1}$ taken in this order in the
  convolution, which occur only if $\sigma_{1}=\sigma_{3}=-\sigma_{2}$ and
  take values in
  $\cH^{6}(M_{2},M_{1})\cong H^{6}\bigl(A,\cO_{A}(-(t_{2}-t_{1})\theta)\bigr)$,
  of dimension $(t_{2}-t_{1})^{6}$.
\item $\dim\Ext^{2}(E,E)\ge249-\min\bigl(249,(t_{2}-t_{1})^{6}\bigr)$, and
  $\dim\Ext^{2}(E,E)\ge249$ unless the order and the signs are as in (i).
  In particular $E$ does not meet the polarised criterion when
  $t_{2}-t_{1}\le2$, since $249-64=185>57$, nor when the signs or the order
  differ from (i).
\item If $E$ meets the polarised criterion, then the length four products
  of (i) have rank at least $249-57=192$ on the kernel of the cup products,
  and $t_{2}-t_{1}\ge3$.
\end{enumerate}
The same holds for $t_{1}<t_{2}<t_{3}\le p-2$ with $E$ replaced by its dual.
\end{theorem}

\begin{proof}
(i) By \Cref{lem:harmonicmodel}(iii) a term with $k\ge3$ leaves acting on
$y\in H^{2}(\cO_{A})$ at a piece $c$ runs along a chain
$s=a_{0}\to\dots\to a_{k-1}=e$ of distinct pieces, increasing in the order
of the convolution, whose $k-1$ steps are components of degree one, with $y$
inserted at $c$. Its value lies in $\cH(s,e)$, in $\Ext$ degree
$2+\sum_{j}q_{j}-(k-2)=3+\varepsilon$, where $\varepsilon$ is the sum of the
excesses of the steps; so $-3\le\varepsilon\le3$, and $\cH(s,e)$ must be
nonzero in degree $3+\varepsilon$. By the proof of \Cref{prop:nothingenters}
a step between two $L_{\zeta}$ has excess at least $1$, with equality only
between pieces of opposite parities, and a step between pieces at different
values of $\phi$ has excess $-1$ and $\Ext$ degree $0$ going up, excess $5$
and $\Ext$ degree $6$ going down. Let $U$ and $D$ be the numbers of steps
up and down and $\lambda$ the total excess of the steps between two
$L_{\zeta}$, so that $\varepsilon=5D-U+\lambda$. For two distinct
$L_{\zeta}$, $L_{\zeta'}$ differing in $m\ge1$ coordinates, $\cH(L_{\zeta},
L_{\zeta'})$ vanishes above degree $6-m$; for pieces at different values it
sits in degree $0$ if the target is higher and in degree $6$ if it is lower.

Two observations remove chains. First, a chain of total excess $-2$ has no
step down, and so consists of steps between $L_{\zeta}$ followed by at most
three steps up. Second, let $a\to b\to c$ be two consecutive steps of
$\Ext$ degree $0$, sections $x_{ab}$ and $x_{bc}$ of line bundles. The
Maurer--Cartan equation $\pi(\delta\circ\delta)=0$ of the proof of
\Cref{prop:explicitconvolution}(i) has, in $\cH^{0}(a,c)$, one term for each
chain from $a$ to $c$ of total excess $-2$. If $a\to b\to c$ is the only such
chain, the equation reads $x_{ab}x_{bc}=0$, a product of holomorphic
sections being holomorphic and so harmonic; since $A$ is connected, $x_{ab}$
or $x_{bc}$ vanishes, and every term along a chain through both steps is
zero.

Here every step up ends at one of $M_{1},M_{2},M_{3}$, so $U\le3$, and
$D\le1$. If $D=0$, the chain is $r\ge0$ pieces $L_{\zeta}$ followed by
multiples in increasing order. When it ends at an $L_{\zeta}$, $k\ge3$ gives
$\lambda\ge2$, so $s$ and $e$ differ in one coordinate and $\lambda=2$; then
the two steps change the parity twice, whereas $s$ and $e$ have opposite
parities. When it ends at a multiple, $e$ lies above $s$, so $\varepsilon=-3$,
which forces $U=3$, $\lambda=0$ and the chain $L_{\zeta}\to M_{1}\to
M_{2}\to M_{3}$; by the first observation $L_{\zeta}\to M_{1}\to M_{2}$ is the
only chain of excess $-2$ from $L_{\zeta}$ to $M_{2}$, and the second removes
it. If $D=1$, then $U\ge2+\lambda$, so $(U,\lambda)$ is $(2,0)$, $(3,0)$ or
$(3,1)$ and $\varepsilon$ is $3$, $2$ or $3$. For $\varepsilon=3$ the piece
$s$ is a multiple and $e$ lies below it. With $(3,1)$ the three steps up
would end at three multiples other than $s$, which do not exist. With $(2,0)$
the chain has three steps, and a step down lands at an $L_{\zeta}$, since a
fourth multiple does not exist. It is $M_{1}\to M_{2}\to M_{3}\to L_{\zeta}$,
removed by the second observation at $M_{1}\to M_{2}\to M_{3}$; or
$M_{a}\to M_{b}\to L_{\zeta}\to M_{c}$ with $t_{a}<t_{b}$ and $t_{c}<t_{a}$,
that is $M_{2}\to M_{3}\to L_{\zeta}\to M_{1}$; or
$M_{a}\to L_{\zeta}\to M_{b}\to M_{c}$ with $t_{b}<t_{c}<t_{a}$, that is
$M_{3}\to L_{\zeta}\to M_{1}\to M_{2}$, removed at $L_{\zeta}\to M_{1}\to
M_{2}$. For $\varepsilon=2$, with $(U,\lambda)=(3,0)$, the pieces $s$ and $e$
are $L_{\zeta}$ differing in one coordinate; the last step is the step down,
preceded by three steps up from $s$, so the chain is $L_{\zeta}\to M_{1}\to
M_{2}\to M_{3}\to L_{\zeta'}$, removed at $L_{\zeta}\to M_{1}\to M_{2}$.

Only $M_{2}\to M_{3}\to L_{\zeta}\to M_{1}$ remains. Its steps have $\Ext$
degrees $0$, $6$ and $0$, so degree one forces $d_{M_{3}}=d_{M_{2}}+1$,
$d_{L}=d_{M_{3}}-5$ and $d_{M_{1}}=d_{L}+1$; this gives the signs, and the
value lies in
$\cH^{3+d_{M_{2}}-d_{M_{1}}}(M_{2},M_{1})=\cH^{6}(M_{2},M_{1})$, whose
dimension is $|\chi(\cO_{A}((t_{1}-t_{2})\theta))|=(t_{2}-t_{1})^{6}$.
(ii) By \Cref{prop:nothingenters} it suffices to bound the kernel of $d_{E}$
on $D^{2}$. By \Cref{prop:cupkernel} the cup products have a kernel $Y$ of
dimension at least $249$, and on $Y$ the differential is the sum of the
products of (i), of rank at most $(t_{2}-t_{1})^{6}$. (iii) follows from
(ii). For the dual, $E\mapsto E^{\vee}$ preserves $\Ext^{2}$ and sends the
configuration with $p$, $t_{i}$ to the one with $-p$, $-t_{i}$, the set of
the $L_{\zeta}$ being stable under $\zeta\mapsto-\zeta$.
\end{proof}

\begin{figure}[tbp]
\centering
  \includegraphics[width=\textwidth]{fig_diagonalbudget}
\caption{The $525$ diagonal classes of degree two at $n=3$
(\Cref{prop:cupkernel,thm:esixdiagonal}). The cup products remove $276$ of
them for general components, and at most that many in any case, leaving
$204$ at the pieces $L_{\zeta}$ and the $45$ at the $M_{i}$; when the $t_{i}$
lie on one side of $p$, the fourfold products of the surviving path remove
at most $(t_{2}-t_{1})^{6}$ more,
which is $1$ or $64$ for the gaps $1$ and $2$ and no constraint from the
third gap on. The criterion asks for $57$, so $468$ classes must disappear
in all and $192$ after the cup products. On the right, the graphs
$\Gamma_{\tau}$ on the $32$ pieces, one row for each value of $\zeta_{1}$:
for $\tau=(2,0,0)$ each row is a component, for $\tau=(1,1,0)$ the
components are the pairs differing in the sign of $\zeta_{3}$; a class
constant on the components of every $\Gamma_{\tau}$ has vanishing cup
products.}
\label{fig:diagonalbudget}
\end{figure}

The placements of the $t_{i}$ on both sides of $p$ are not covered by
\Cref{thm:esixdiagonal}. There two families of fourfold products occur
instead of one, and no convolution carries both.

\begin{proposition}[Placements on both sides of $p$]\label{prop:mixedplacement}
Keep \Cref{thm:esixdiagonal}, but let $t_{1}\le p-2$ and
$p+2\le t_{2}<t_{3}$.
\begin{enumerate}[label=\textup{(\roman*)},leftmargin=2.6em,itemsep=2pt]
\item The Massey products of length at least three vanish on $D^{2}$ unless
  $\sigma_{1}=\sigma_{2}=-\sigma_{3}$. In that case the only ones that act
  have length four and run along one of the paths
  \[
    \textup{(a)}\ \ M_{2}\to M_{3}\to M_{1}\to L_{\zeta},
    \qquad
    \textup{(b)}\ \ M_{3}\to M_{1}\to L_{\zeta}\to M_{2},
  \]
  taken in this order in the convolution, with $\prod\zeta=\sigma_{3}$. The
  products along (a) take values in $\cH^{6}(M_{2},L_{\zeta})$, of dimension
  $\bigl((t_{2}-p)^{2}-1\bigr)^{3}$, and those along (b) in
  $\cH^{6}(M_{3},M_{2})$, of dimension $(t_{3}-t_{2})^{6}$. A convolution
  carries at most one of the two kinds, since (a) puts $M_{2}$ before $M_{3}$
  and (b) puts $M_{3}$ before $M_{2}$.
\item If $E$ meets the polarised criterion, then
  $\sigma_{1}=\sigma_{2}=-\sigma_{3}$, and the products along (a) or those
  along (b) have rank at least $192$ on the kernel of the cup products. For
  (b) this needs $t_{3}-t_{2}\ge3$. The sixteen targets of (a) have
  dimension $16\bigl((t_{2}-p)^{2}-1\bigr)^{3}\ge432$, so for (a) the
  count gives no condition.
\end{enumerate}
For $t_{1}<t_{2}\le p-2$ and $t_{3}\ge p+2$ the same holds for
$E^{\vee}[1]$, and the condition on the signs becomes
$\sigma_{2}=\sigma_{3}=-\sigma_{1}$.
\end{proposition}

\begin{proof}
(i) We use the notation and the two observations of the proof of
\Cref{thm:esixdiagonal}(i). Now $\phi(M_{1})=t_{1}<p<t_{2}<t_{3}$; a step up
ends at $M_{2}$, at $M_{3}$, or at an $L_{\zeta}$ coming from $M_{1}$, so
again $U\le3$ and $D\le1$. If $D=0$, the values do not decrease along the
chain. When $e$ lies above $s$, $\varepsilon=-3$ forces $U=3$, $\lambda=0$
and the chain $M_{1}\to L_{\zeta}\to M_{2}\to M_{3}$, removed at
$L_{\zeta}\to M_{2}\to M_{3}$, the only chain of excess $-2$ from
$L_{\zeta}$ to $M_{3}$. When $e$ and $s$ have the same value, they are two
$L_{\zeta}$ joined by steps between $L_{\zeta}$, which is excluded as before.
If $D=1$, $(U,\lambda)$ is $(2,0)$, $(3,0)$ or $(3,1)$. For
$\varepsilon=3$, $e$ lies below $s$. With $(3,1)$ the steps up end at
$M_{2}$, $M_{3}$ and an $L_{\zeta}$ following $M_{1}$; then $s$ is an
$L_{\zeta}$, so $e=M_{1}$, which is followed by a step: impossible. With
$(2,0)$ the chain has three steps and no two consecutive $L_{\zeta}$. As
$M_{1}$ is the lowest piece, it is never the end of a step up, and a step up
into an $L_{\zeta}$ comes from $M_{1}$; so two consecutive steps up end at
$M_{2}$ or $M_{3}$. The chain cannot start at $M_{1}$. If it starts at an
$L_{\zeta}$, it ends at $M_{1}$ with the step down, after two steps up:
$L_{\zeta}\to M_{2}\to M_{3}\to M_{1}$, removed at $L_{\zeta}\to M_{2}\to
M_{3}$. If it starts at $M_{3}$, the first step goes down and the last two
go up, ending at $M_{2}$; the middle piece is then an $L_{\zeta}$ entered
from $M_{1}$, and the chain is $M_{3}\to M_{1}\to L_{\zeta}\to M_{2}$, path
(b). If it starts at $M_{2}$, then $e\in\{L_{\zeta},M_{1}\}$, so the last two
steps are not both up; the first step goes up to $M_{3}$, the second down,
and the third up into $e=L_{\zeta}$ from $M_{1}$: the chain is $M_{2}\to
M_{3}\to M_{1}\to L_{\zeta}$, path (a). For $\varepsilon=2$,
$(U,\lambda)=(3,0)$ and $s$, $e$ are $L_{\zeta}$ differing in one
coordinate. The last step enters $e$ from $M_{1}$, since a step down into
$e$ would follow three steps up from $s$, one more than the values allow;
$M_{1}$ is entered by the step down, after two steps up from $s$:
$L_{\zeta}\to M_{2}\to M_{3}\to M_{1}\to L_{\zeta'}$, removed at
$L_{\zeta}\to M_{2}\to M_{3}$.
Along (a) the steps have $\Ext$ degrees $0$, $6$ and $0$, so degree one
forces $d_{M_{3}}=d_{M_{2}}+1$, $d_{M_{1}}=d_{M_{3}}-5$ and
$d_{L}=d_{M_{1}}+1$. Along (b) they have $\Ext$ degrees $6$, $0$ and $0$,
which forces $d_{M_{1}}=d_{M_{3}}-5$, $d_{L}=d_{M_{1}}+1$ and
$d_{M_{2}}=d_{L}+1$. This gives the signs and puts the values in
$\cH^{3+d_{M_{2}}-d_{L}}(M_{2},L_{\zeta})=\cH^{6}(M_{2},L_{\zeta})$ and in
$\cH^{3+d_{M_{3}}-d_{M_{2}}}(M_{3},M_{2})=\cH^{6}(M_{3},M_{2})$. On $B_{j}$
the line bundle $L_{\zeta}\otimes M_{2}^{-1}$ has the hermitian form
$\bigl(\begin{smallmatrix}p-t_{2}&\zeta_{j}\\ \bbar\zeta_{j}&p-t_{2}
\end{smallmatrix}\bigr)$, with eigenvalues $p-t_{2}\pm1<0$ and determinant
$(t_{2}-p)^{2}-1$, which gives the first dimension. The second is that of
\Cref{thm:esixdiagonal}(i) with $t_{3}$, $t_{2}$ in place of $t_{2}$,
$t_{1}$.
(ii) \Cref{prop:nothingenters,prop:cupkernel} do not depend on the
placement of the $t_{i}$, so as in \Cref{thm:esixdiagonal}(ii)
$\dim\Ext^{2}(E,E)\ge249-\rho$, where $\rho$ is the rank on the kernel of
the cup products of the products of (i) that $E$ carries, and $\rho=0$
unless $\sigma_{1}=\sigma_{2}=-\sigma_{3}$. The criterion asks for $57$,
so $\rho\ge192$, and $(t_{3}-t_{2})^{6}\ge192$ exactly when
$t_{3}-t_{2}\ge3$. For the last statement, $E\mapsto E^{\vee}[1]$ preserves
$\Ext^{2}$, sends $p$ and the $t_{i}$ to $-p$ and the $-t_{i}$, reverses the
order of the pieces and changes every sign.
\end{proof}

The Maurer--Cartan equation can be solved for one order of the pieces.

\begin{proposition}[An explicit convolution]\label{prop:explicitconvolution}
Keep \Cref{thm:esixdiagonal} with $p=0$ and $t=(2,5,6)$, take the pieces to
be external products $L_{\zeta}=N_{1,\zeta_{1}}\boxtimes N_{2,\zeta_{2}}
\boxtimes N_{3,\zeta_{3}}$ and $M_{i}=M_{i,1}\boxtimes M_{i,2}\boxtimes
M_{i,3}$ of line bundles on the $B_{j}$, and order them as $M_{2}$, $M_{3}$,
the sixteen $L_{\zeta}$ with $\prod\zeta=1$, $M_{1}$, the sixteen $L_{\zeta}$
with $\prod\zeta=-1$, with the shifts $0$, $1$, $-4$, $-3$, $-5$.
\begin{enumerate}[label=\textup{(\roman*)},leftmargin=2.6em,itemsep=2pt]
\item The components that the degrees allow are
  $x_{1}\in\cH^{0}(M_{2},M_{3})$, $x_{2}^{\zeta}\in\cH^{6}(M_{3},L_{\zeta})$,
  $x_{3}^{\zeta}\in\cH^{0}(L_{\zeta},M_{1})$, components in
  $\cH^{2}(L_{\zeta},L_{\zeta'})$ for the $48+96$ pairs with
  $\prod\zeta=1=-\prod\zeta'$ that differ in one or two coordinates, and
  components in $\cH^{6}(M_{2},L_{\zeta'})$, $\prod\zeta'=-1$, which we take
  to be zero. For harmonic components the Maurer--Cartan equation is exactly
  \[
    x_{1}x_{2}^{\zeta}=0\ \text{in}\ H^{6}(A,M_{2}^{-1}\otimes L_{\zeta})
    \ \text{for every}\ \zeta,
    \qquad
    \sum_{\zeta}x_{2}^{\zeta}x_{3}^{\zeta}=0\ \text{in}\
    H^{6}(A,M_{3}^{-1}\otimes M_{1}),
  \]
  and every solution is the harmonic part of a convolution $E$ with
  $\ch(E)=32(\alpha+\bbar\alpha)+e^{5\theta}-e^{6\theta}-e^{2\theta}$.
\item There are solutions with every $x_{2}^{\zeta}\ne0$ and
  $x_{1}x_{2}^{\zeta}=x_{2}^{\zeta}x_{3}^{\zeta}=0$.
\item For every solution $\dim\Ext^{2}(E,E)\ge2560$, so $E$ does not meet
  the polarised criterion.
\end{enumerate}
\end{proposition}

\begin{proof}
(i) A component from $P_{a}$ to $P_{b}$ has $\Ext$ degree $1+d_{a}-d_{b}$, and
the forms give the cohomology of each pair: $\cH(M_{2},M_{3})$ and
$\cH(L_{\zeta},M_{1})$ sit in degree $0$, $\cH(M_{3},L_{\zeta})$,
$\cH(M_{2},L_{\zeta'})$ and $\cH(M_{2},M_{1})$ in degree $6$, and
$\cH(L_{\zeta},L_{\zeta'})$ in degree $2$ exactly for the pairs listed; no
other pair has cohomology in the degree a component needs. Put
$\delta=\delta^{H}-h(\delta\circ\delta)$ recursively, as in
\Cref{lem:harmonicmodel}(i), and $\Phi=\bar\partial\delta+\delta\circ\delta$.
Then $\Phi=\pi(\delta\circ\delta)+h\bar\partial(\delta\circ\delta)$ and
$\bar\partial(\delta\circ\delta)=\Phi\circ\delta-\delta\circ\Phi$, so by
induction on $b-a$ the twisted complex exists exactly when
$\pi(\delta\circ\delta)=0$. The components of $\pi(\delta\circ\delta)$ come
from paths of two or more components; those that land in a nonzero group
are $x_{1}x_{2}^{\zeta}$ in $\cH^{6}(M_{2},L_{\zeta})$ and
$\sum_{\zeta}x_{2}^{\zeta}x_{3}^{\zeta}$ in $\cH^{6}(M_{3},M_{1})$, while the
product along $M_{2}\to M_{3}\to L_{\zeta}\to M_{1}$ would lie in
$\cH^{5}(M_{2},M_{1})=0$ and $x_{2}^{\zeta}$ composed with a component in
$\cH^{2}(L_{\zeta},L_{\zeta'})$ in degree eight. The signs
$(-1)^{d}$ are $+,-,+,-,-$, which gives the Chern character by
\Cref{prop:weilpieces}(i).
(ii) Take $x_{1}$ and the $x_{3}^{\zeta}$ to be external products of
nonzero sections on the $B_{j}$. On each $B_{j}$ the products with their
factors map $H^{2}(B_{j},M_{3}^{-1}L_{\zeta})$, of dimension $35$, to spaces
of dimensions $24$ and $16$, so they have kernels of dimensions at least $11$
and $19$. A sum of external products whose factor on one surface lies in the
first kernel and whose factor on another lies in the second is a nonzero
$x_{2}^{\zeta}$ with $x_{1}x_{2}^{\zeta}=x_{2}^{\zeta}x_{3}^{\zeta}=0$, which
solves (i).\footnote{Such solutions are far from unique: the equations of
(i) are linear in each of $x_{1}$, $x_{2}^{\zeta}$, $x_{3}^{\zeta}$
separately, and the components on the $144$ pairs do not enter them at
all.}
(iii) The sixteen $L_{\zeta}$ with $\prod\zeta=1$ have shift $-4$ and the
others $-5$, so \Cref{lem:allthree} below applies to every pair of them
that differ in all three coordinates: there are $112$ such pairs, and their
groups $\cH^{3}$ have dimensions adding up to $16\cdot160=2560$.
\end{proof}

The classes counted in (iii) are the ones no product can reach. Two pieces
$L_{\zeta}$, $L_{\zeta'}$ that differ in every coordinate have a difference
that is indefinite on every surface, so $\cH(L_{\zeta},L_{\zeta'})$ is
$\bigotimes_{j}H^{1}(B_{j},\cdot)$, concentrated in degree three, of
dimension $D(\zeta,\zeta')=\prod_{j}|\zeta_{j}-\zeta'_{j}|^{2}$; no component
can join them, and if their shifts differ by one their classes have degree
two.

\begin{lemma}[Pieces that differ everywhere]\label{lem:allthree}
Keep \Cref{prop:nothingenters} with $n=3$. Suppose that for some integer $r$
either $p+2\le t_{1}<t_{2}<t_{3}$ and $M_{1}$, $M_{2}$, $M_{3}$ have the
shifts $r+1$, $r+4$, $r+5$, or $t_{1}\le p-2<p+2\le t_{2}<t_{3}$ and they
have the shifts $r-1$, $r+3$, $r+4$ or $r-1$, $r+1$, $r+4$; these are the
shifts forced by the paths of \Cref{thm:esixdiagonal}(i) and of
\Cref{prop:mixedplacement}(i), (a) and (b), with $r$ the shift of the
$L_{\zeta}$ on the path. Let $L_{\zeta}$, $L_{\zeta'}$ differ in all three
coordinates, with $d_{L_{\zeta}}=d_{L_{\zeta'}}+1$. Then $d_{E}$ vanishes on
$\cH^{3}(L_{\zeta},L_{\zeta'})$, and no element of degree one has a
differential with a nonzero component there. Consequently
\[
  \dim\Ext^{2}(E,E)\;\ge\;\dim\ker\bigl(d_{E}|_{D^{2}}\bigr)
  +\sum D(\zeta,\zeta'),
\]
the sum running over the pairs of this kind.
\end{lemma}

\begin{proof}
Only the shifts are used, and the order of the pieces only removes
components, so we ignore it. A component from $P_{a}$ to $P_{b}$ of $\Ext$
degree $q$ has $d_{b}=d_{a}+1-q$. Between distinct $L_{\zeta}$ the degree $q$
is at least $2$: there is no $\cH^{0}$, and $\cH^{1}$ occurs only for pieces
differing in one coordinate, which have opposite parities, while $q=1$ would
need equal shifts. So a component between two $L_{\zeta}$ lowers the shift
by at least one, and by exactly one only in degree two. For a multiple
$M_{i}$ let $\mu_{i}$ and $\nu_{i}$ be the shifts of an $L_{\zeta}$ that a
component from $M_{i}$ reaches and from which a component reaches $M_{i}$:
$\mu_{i}=d_{M_{i}}-5$ and $\nu_{i}=d_{M_{i}}-1$ when $t_{i}>p$, and
$\mu_{i}=d_{M_{i}}+1$, $\nu_{i}=d_{M_{i}}+5$ when $t_{i}<p$. Between the
multiples the degrees allow only $M_{2}\to M_{3}$ in the first case,
$M_{2}\to M_{3}\to M_{1}$ in the second, and $M_{3}\to M_{1}$ in the third.
Let $\lambda_{i}$ be the largest $\mu_{k}$ and $\kappa_{i}$ the least
$\nu_{k}$ over the multiples $M_{k}$ reachable from $M_{i}$, respectively
reaching $M_{i}$, through these steps, $M_{i}$ included. The values, in the
order $M_{1}$, $M_{2}$, $M_{3}$ and relative to $r$, are
\[
\begin{array}{l|ccc|ccc|ccc|ccc}
 & \multicolumn{3}{c|}{\mu} & \multicolumn{3}{c|}{\nu}
 & \multicolumn{3}{c|}{\lambda} & \multicolumn{3}{c}{\kappa}\\
\hline
r+1,\,r+4,\,r+5 & -4&-1&0 & 0&3&4 & -4&0&0 & 0&3&3\\
r-1,\,r+3,\,r+4 & 0&-2&-1 & 4&2&3 & 0&0&0 & 2&2&2\\
r-1,\,r+1,\,r+4 & 0&-4&-1 & 4&0&3 & 0&-4&0 & 3&0&3
\end{array}
\]
and in all three cases, for all $i$ and $k$,
\begin{equation}\label{eq:allthree}
  \lambda_{i}\le r\le\kappa_{k},\qquad
  \kappa_{i}\ge\lambda_{i}+2,\qquad
  \lambda_{i}\le\mu_{i}+2,\qquad
  \nu_{i}\le\kappa_{i}+2,
\end{equation}
while $\mu_{i}\le\lambda_{i}$ and $\nu_{i}\ge\kappa_{i}$ by definition. A
run $L\to M_{k}\to\dots\to M_{i}\to L'$ through multiples goes from the shift
$\nu_{k}\ge\kappa_{k}$ to $\mu_{i}\le\lambda_{k}\le\kappa_{k}-2$, so every chain
of components from one $L$ to another lowers the shift, by at least two if it
passes through a multiple and by exactly one only if it is a single component
of degree two. In particular a chain from an $L$ of shift $x$ into $M_{i}$
has $x\ge\kappa_{i}$, and a chain from $M_{i}$ to an $L$ of shift $y$ has
$y\le\lambda_{i}$.

By \Cref{lem:harmonicmodel}(iii), a term of $d_{E}$ on a class $u$ of $\Ext$
degree $q_{u}$ in $\cH(a',b')$ runs along a chain $c\to\dots\to a'$, the
step $u$, and a chain $b'\to\dots\to e$, and since the degrees of the
components telescope, its value lies in $\cH(c,e)$ in $\Ext$ degree
\[
  q_{u}+1+(d_{c}-d_{a'})+(d_{b'}-d_{e}).
\]
Write $a=L_{\zeta}$ and $b=L_{\zeta'}$, so that $d_{a}=d_{b}+1$.

\emph{Nothing reaches $\cH^{3}(a,b)$.} Let $u$ have degree one, so
$q_{u}=1+d_{a'}-d_{b'}$, with chains from $a$ to $a'$ and from $b'$ to $b$
and $(a',b')\ne(a,b)$. If $a'$ and $b'$ are both pieces $L$, then
$d_{a'}\le d_{a}$ and $d_{b'}\ge d_{b}$, not both with equality, so
$q_{u}\le1$; $q_{u}=0$ is impossible, and $q_{u}=1$ needs either $a'=b'$, and
then a chain from $a$ to $b$ lowering the shift by one, a single component in
$\cH^{2}(a,b)=0$, or $a'$, $b'$ differing in one coordinate, of opposite
parities, against $d_{a'}=d_{b'}$. If $a'=M_{i}$ and $b'$ is an $L$, then
$u$ lies in the degree of a component, so $d_{b'}=\mu_{i}$, while
$d_{b'}\ge d_{b}=d_{a}-1\ge\kappa_{i}-1\ge\lambda_{i}+1>\mu_{i}$. If $a'$ is
an $L$ and $b'=M_{k}$, then $d_{a'}=\nu_{k}$ and
$d_{a}-d_{b}\ge\nu_{k}-\lambda_{k}\ge\kappa_{k}-\lambda_{k}\ge2$. If
$a'=M_{i}$ and $b'=M_{k}$, then $d_{a}-d_{b}\ge\kappa_{i}-\lambda_{k}$; for
$i=k$ this is at least $2$, and for $i\ne k$ the class $u$ lies in the
degree of a step $M_{i}\to M_{k}$, so $\lambda_{k}\le\lambda_{i}$ and again
$d_{a}-d_{b}\ge2$. Each case contradicts $d_{a}-d_{b}=1$.

\emph{$d_{E}$ vanishes on $\cH^{3}(a,b)$.} Here $u$ is the class itself and a
term lands in $\cH(c,e)$, $(c,e)\ne(a,b)$, in degree $3+d_{c}-d_{e}$, with
chains from $c$ to $a$ and from $b$ to $e$. If $c$ and $e$ are both pieces
$L$, then $d_{c}-d_{e}\ge2$; $c=e$ is impossible, and otherwise a class of
degree at least $5$ needs $c$, $e$ differing in one coordinate, in degree
$5$, so $d_{c}-d_{e}=2$ against their opposite parities. If $c=M_{i}$ and
$e$ is an $L$, the degree is that of a component when $d_{e}=\mu_{i}+2$, but
$d_{e}\le d_{b}<d_{a}\le\lambda_{i}\le\mu_{i}+2$. If $c$ is an $L$ and
$e=M_{k}$, it is that of a component when $d_{c}=\nu_{k}-2$, but
$d_{c}\ge d_{a}>d_{b}\ge\kappa_{k}\ge\nu_{k}-2$. If $c=M_{i}$ and
$e=M_{k}$, then $d_{a}\le\lambda_{i}\le r\le\kappa_{k}\le d_{b}$, against
$d_{a}>d_{b}$.

For the inequality: a coboundary has no component on $D^{2}$
(\Cref{prop:nothingenters}) and none on these groups, so the cocycles of
$\ker(d_{E}|_{D^{2}})$ and of these groups span a space that meets the
coboundaries only in zero. Twisting the pieces by line bundles of degree
zero does not change $D(\zeta,\zeta')$, the forms being nondegenerate on
every surface, and can only remove components.
\end{proof}

\begin{theorem}[No convolution of these pieces meets the criterion]
\label{thm:noconvolution}
Keep \Cref{prop:nothingenters} with $n=3$: let $E$ be any convolution of
the $32$ bundles $L_{\zeta}$ and of $M_{1}$, $M_{2}$, $M_{3}$, with
$t_{1}<t_{2}<t_{3}$ and $|t_{i}-p|\ge2$, in any order, with any shifts of the
given parities and with the pieces twisted by any line bundles of degree
zero. Then $\dim\Ext^{2}(E,E)\ge69$. In particular no such convolution meets
the polarised criterion, which asks for $57$.
\end{theorem}

\begin{proof}
Replacing $E$ by $E^{\vee}$ or $E^{\vee}[1]$, which preserves $\Ext^{2}$, we
may assume that the $t_{i}$ are as in \Cref{thm:esixdiagonal} or in
\Cref{prop:mixedplacement}. If no fourfold product of those results acts,
$\dim\Ext^{2}(E,E)\ge249$. Otherwise the path fixes the shifts of the
$M_{i}$ as in \Cref{lem:allthree}, with $r$ the shift of every $L_{\zeta}$
on a path; let $P$ be the sixteen pieces of the parity of $r$, $Q$ the
other sixteen, and $m$ the number of pieces on paths.

The fourfold products act on $D^{2}$ only through the classes at the
$M_{i}$ and at the $m$ pieces on paths, so their rank is at most $45+15m$,
and by \Cref{prop:cupkernel}
$\dim\ker(d_{E}|_{D^{2}})\ge249-45-15m=204-15m$. If $m\le9$ this is at
least $69$.

Let $m\ge10$, so that at least ten pieces of $P$ have the shift $r$. Every
$W\in Q$ has $3$, $6$ and $7$ partners in $P$ differing from it in one, two
and three coordinates, and for the seven $D=16$ six times and $D=64$
once.\footnote{The ratios $\zeta'_{j}/\zeta_{j}$ of two pieces of opposite
parities have product $-1$; when all three differ from $1$ they are
$(-1,-1,-1)$, with $D=4^{3}$, or a permutation of $(i,-i,-1)$, with
$D=2\cdot2\cdot4$.} If $d_{W}=r\pm1$, at most six pieces of $P$ lack the
shift $r$, so one of the seven has it, and \Cref{lem:allthree} gives at least
$16$ classes. If $d_{W}\ne r\pm1$, consider for each type $\tau$ of the form
$(1,1,0)$ the partner $W^{\tau}\in P$ obtained by changing the sign of the
coordinate where $\tau$ vanishes. By the proof of \Cref{prop:cupkernel},
$W^{\tau}$ is the only piece with which $W$ can carry a nonzero cup product
on the type $\tau$, and conversely; a cup product between them needs a
component of degree two, hence shifts differing by one. If they do, the
classes of type $\tau$ equal at $W$ and $W^{\tau}$ and zero elsewhere have
vanishing cup products, and $W^{\tau}$, of shift $d_{W}\pm1\ne r$, lies on no
path; if they do not, the classes of type $\tau$ at $W$ alone have vanishing
cup products. Either way these $4$ classes lie in $\ker(d_{E}|_{D^{2}})$,
$12$ for each such $W$, and those of different $W$ are independent, since
$W\mapsto W^{\tau}$ is injective. Adding over $Q$, \Cref{lem:allthree}
gives $\dim\Ext^{2}(E,E)\ge16\cdot12=192$.
\end{proof}

At $n=4$ there are $128$ pieces $L_{\zeta}$, and the classes that no
product reaches are again those between pieces that differ in three
coordinates and have adjacent shifts. The shifts along a chain are now
controlled by counting alone.

\begin{lemma}[Shifts along chains at $n=4$]\label{lem:efourshifts}
Keep \Cref{prop:weilpieces} with $n=4$, take the $L_{\zeta}$ to be external
products $N_{1,\zeta_{1}}\boxtimes\dots\boxtimes N_{4,\zeta_{4}}$ of line
bundles on the $B_{j}$, let $t_{1},t_{2},t_{3}$ be distinct with
$|t_{i}-p|\ge2$, and let $E$ be a convolution of the $L_{\zeta}$ and of
$M_{1},M_{2},M_{3}$ with shifts of the given parities. Call a component of
$\delta^{H}$, or an element of degree one of $\cH_{E}$ in one group
$\cH(a,b)$, a step from $P_{a}$ to $P_{b}$; a step of $\Ext$ degree $q$ has
$d_{b}=d_{a}+1-q$. Put $\phi(L_{\zeta})=p$ and $\phi(M_{i})=t_{i}$.
\begin{enumerate}[label=\textup{(\roman*)},leftmargin=2.6em,itemsep=2pt]
\item A step between distinct $L_{\zeta}$ has $\Ext$ degree at least two, so
  it lowers the shift by at least one, and by exactly one only in $\Ext$
  degree two. A step in $\cH^{1}(a,a)=H^{1}(A,\cO_{A})$ keeps the shift.
\item A step between an $L_{\zeta}$ and a multiple, or between two
  multiples, has $\Ext$ degree $0$ if it goes up the value $\phi$ and $8$ if
  it goes down, so it raises the shift by one or lowers it by seven.
\item A sequence of steps from an $L_{\zeta}$ to an $L_{\zeta'}$ through
  multiples, distinct before one marked step and distinct after it, where
  the marked step may lie in $\cH^{1}(a,a)$ at a multiple, lowers the shift
  by at least four. Consequently a sequence of steps from one $L_{\zeta}$ to
  another, at most one of them not a component, lowers the shift by at least
  one, and by exactly one only if it consists of steps between pieces
  $L_{\zeta}$, one of $\Ext$ degree two and the others in $\cH^{1}(a,a)$.
\end{enumerate}
\end{lemma}

\begin{proof}
(i) For $\zeta\ne\zeta'$ the group $\cH(L_{\zeta},L_{\zeta'})$ is the tensor
product over the surfaces of $H^{1}$ of a line bundle with a nondegenerate
form, on the $k\ge1$ surfaces where $\zeta$ and $\zeta'$ differ, and of
$H^{*}(B_{j},\cO_{B_{j}})$ on the others, the factors of $L_{\zeta}$ and
$L_{\zeta'}$ on those being equal. So its degrees are at least $k$, and degree
one needs $k=1$. Then the ratio $\zeta'_{j}/\zeta_{j}\ne1$ in the one
coordinate is the ratio of $\prod\zeta'$ and $\prod\zeta$, which is $-1$, and
the shifts have opposite parities, while a step of $\Ext$ degree one keeps
the shift. (ii) The forms of $M_{i}\otimes L_{\zeta}^{-1}$ have the
eigenvalues $t_{i}-p\pm1$ on every surface, which have the sign of $t_{i}-p$;
the form of $M_{k}\otimes M_{i}^{-1}$ is $(t_{k}-t_{i})\,\mathrm{Id}$. A
definite form has cohomology in degree $0$ if it is positive and $2n=8$ if
it is negative \cite[Chapter~3]{BL04}.%
\footnote{This is Mumford's index theorem: a line bundle on an abelian
variety whose hermitian form is nondegenerate with $i$ negative eigenvalues
has cohomology only in degree $i$. On $B_{j}=E_{0}\times E_{0}$ the form of
$M_{i}\otimes L_{\zeta}^{-1}$ is
$\bigl(\begin{smallmatrix}t_{i}-p&-\zeta_{j}\\-\bar\zeta_{j}&t_{i}-p\end{smallmatrix}\bigr)$,
with eigenvalues $t_{i}-p\pm1$, and $|t_{i}-p|\ge2$ makes both of the sign of
$t_{i}-p$; over the four surfaces the index is $0$ or $8$. The hypothesis
$|t_{i}-p|\ge2$ is what keeps the steps between a piece and a multiple out
of the middle degrees.}
(iii) Let $U$ and $D$ be the numbers of steps up and down outside
$\cH^{1}(a,a)$; by (ii) the shift changes by $U-7D$. The values start and
end at $p$ and pass through the values of multiples, so $D\ge1$. If $D=1$,
the values increase strictly before the step down and strictly after it, so
those before it lie above $p$ and those after it below $p$, each multiple
occurs once, and $U\le3$: the drop $7-U$ is at least four.%
\footnote{The bound is attained: with $p$ below the three values $t_{i}$, the
run $L_{\zeta}\to M_{1}\to M_{2}\to M_{3}\to L_{\zeta'}$, up three times and
down once, changes the shift by $3-7=-4$. With $p$ in the middle,
$t_{1}<p<t_{2}<t_{3}$, the run up through $M_{2}$ and $M_{3}$, down to $M_{1}$
and up to $L_{\zeta'}$ changes it by $3-7=-4$ as well.} If $D\ge2$, at most
six multiples occur, counted with repetition, so there are at most seven
steps outside $\cH^{1}(a,a)$, $U\le7-D$, and the drop $7D-U\ge8D-7$ is at
least nine. For the consequence, split the sequence at the pieces
$L_{\zeta}$ it visits: each part through multiples lowers the shift by at
least four, and by (i) each other part lowers it by at least one, unless it
lies in $\cH^{1}(a,a)$.
\end{proof}

\begin{theorem}[Two shifts on $E_{0}^{8}$]\label{thm:efourtwolevels}
Keep \Cref{lem:efourshifts}, and let $L_{\zeta}$ and $L_{\zeta'}$ differ in
exactly three coordinates and have the shifts $d_{L_{\zeta}}=d_{L_{\zeta'}}+1$.
Then $\cH(L_{\zeta},L_{\zeta'})$ has its classes of degree two in
$\cH^{3}(L_{\zeta},L_{\zeta'})$, of dimension
$D(\zeta,\zeta')=\prod_{\zeta_{j}\ne\zeta'_{j}}|\zeta_{j}-\zeta'_{j}|^{2}$,
which is $16$ or $64$, and:
\begin{enumerate}[label=\textup{(\roman*)},leftmargin=2.6em,itemsep=2pt]
\item for every convolution, no element of degree one has a differential
  with a nonzero component in $\cH^{3}(L_{\zeta},L_{\zeta'})$;
\item if the $L_{\zeta}$ occupy two consecutive shifts, those with
  $\prod\zeta=\epsilon$ the shift $s$ and the others $s-1$, for a sign
  $\epsilon$ and an integer $s$, then $d_{E}$ vanishes on these groups, and
  for every order of the pieces, every choice of the shifts of the $M_{i}$
  and every placement of the $t_{i}$,
  \[
    \dim\Ext^{2}(E,E)\;\ge\;64\cdot640=40960 .
  \]
  In particular $E$ does not meet the polarised criterion, which asks for
  $\dim\Ext^{2}(E,E)=104$.
\end{enumerate}
\end{theorem}

\begin{proof}
Write $W=L_{\zeta}$ and $V=L_{\zeta'}$. By the proof of
\Cref{lem:efourshifts}(i), $\cH(W,V)$ has degrees $3$, $4$ and $5$, the last
two from the surface where $\zeta$ and $\zeta'$ agree, and its classes of
degree $w=2$ have $\Ext$ degree $2+d_{W}-d_{V}=3$; the dimension of
$\cH^{3}(W,V)$ is the product over the three other surfaces of the
determinants $|\zeta_{j}-\zeta'_{j}|^{2}$, as in \Cref{lem:allthree}.

(i) As in the proof of \Cref{lem:allthree}, a term of the differential of an
element $u$ of degree one in $\cH(a',b')$ with a component in $\cH(W,V)$
runs along a chain of components $W\to\dots\to a'$, the step $u$, and a chain
$b'\to\dots\to V$, with $(a',b')\ne(W,V)$. This is a sequence of steps from
$W$ to $V$, one of them $u$, lowering the shift by one. By
\Cref{lem:efourshifts}(iii) it consists of steps between pieces $L_{\zeta}$,
one of $\Ext$ degree two and at most one, $u$, in $\cH^{1}(a,a)$. If $u$ is
not in $\cH^{1}(a,a)$, it is the step of degree two and both chains are
empty, against $(a',b')\ne(W,V)$. If $u$ is in $\cH^{1}(a,a)$, the rest is a
single component $W\to V$ of $\Ext$ degree two, in $\cH^{2}(W,V)=0$.

(ii) A term of $d_{E}$ on $x\in\cH^{3}(W,V)$ runs along chains of components
$c\to\dots\to W$ and $V\to\dots\to e$, not both empty, and lands in $\cH(c,e)$
in $\Ext$ degree $3+1+(d_{c}-d_{W})+(d_{V}-d_{e})=3+d_{c}-d_{e}$. No piece
$L_{\zeta''}$ other than $W$ lies on the first chain: by
\Cref{lem:efourshifts}(iii) it would have the shift at least $d_{W}+1=s+1$.
Likewise only multiples follow $V$ on the second chain, and each chain has at
most three pieces besides $W$ or $V$. Let $U$ and $D$ be the numbers of steps
up and down along both chains, so that $U+D\le6$; by
\Cref{lem:efourshifts}(ii), $d_{W}=d_{c}+U_{1}-7D_{1}$ and
$d_{e}=d_{V}+U_{2}-7D_{2}$, and the $\Ext$ degree is
\[
  3+d_{c}-d_{e}=4-U+7D .
\]
If $c=e$, this degree is $3$, so $U=7D+1$, and $D\ge1$ because the values
return to $\phi(c)\ne p$; then $U\ge8$, which is impossible. If $c\ne e$, the
group $\cH(c,e)$ sits in degree $0$ if $\phi(e)>\phi(c)$ and in degree $8$ if
$\phi(e)<\phi(c)$, by \Cref{lem:efourshifts}(ii). Degree $0$ needs
$U=4+7D\ge4$ with $D=0$, so the values increase strictly, the multiples are
distinct and $U\le3$. Degree $8$ needs $7D=U+4$, so $D=1$ and $U=3$; the
values increase strictly before the step down, from $\phi(c)$, and after it,
up to $\phi(e)<\phi(c)$, so the five values along the path, with $W$ and $V$
counted once, are distinct, and four of them are values of multiples, of
which there are three. So no term reaches a nonzero group, and $d_{E}$
vanishes on $\cH^{3}(W,V)$.

By (i) and (ii) the classes of these groups are cocycles of degree two, and
no nonzero combination of them is a coboundary, since every coboundary has
zero components there. Each piece of the shift $s$ has $28$ partners of the
other parity differing from it in exactly three
coordinates,\footnote{The ratios $\zeta'_{j}/\zeta_{j}$ in the three
coordinates have product $-1$, so they are $(-1,-1,-1)$, with $D=64$, or a
permutation of $(i,-i,-1)$, with $D=16$; with four places for the
coordinate where they agree this gives $4\cdot7=28$ partners and
$4\cdot(64+6\cdot16)=640$ classes.} all of the shift $s-1$, and the $28$
groups have dimensions adding up to $640$; there are $64$ pieces of the shift
$s$.
\end{proof}

\begin{figure}[tbp]
\centering
\includegraphics[width=\textwidth]{fig_efourshifts}
\caption{\Cref{lem:efourshifts} and \Cref{thm:efourtwolevels}, in the
placement $p<t_{1}<t_{2}<t_{3}$. Horizontal: the value $\phi$ of a piece;
vertical: its shift. The pieces $L_{\zeta}$ with $\prod\zeta=\epsilon$ sit at
the shift $s$ and the others at $s-1$. A run through the multiples, up three
times with $\Ext$ degree $0$ and down once with $\Ext$ degree $8$, ends four
shifts lower, the least drop of \Cref{lem:efourshifts}(iii), so it never
joins a top piece $W$ to a partner $V$ one shift lower; the group
$\mathcal{H}^{3}(W,V)$ survives in $\Ext^{2}(E,E)$.}
\label{fig:efourshifts}
\end{figure}

\Cref{fig:efourshifts} draws the shifts along a run. The hypothesis of (ii) is the arrangement of the explicit convolution of
\Cref{prop:explicitconvolution} at $n=3$, where the sixteen pieces of each
parity share a shift. The pieces
must not be twisted by line bundles of degree zero here, unlike in
\Cref{thm:noconvolution}: a twist that is nontrivial on the surface where
$\zeta$ and $\zeta'$ agree removes $\cH(L_{\zeta},L_{\zeta'})$.

With a third shift, a component of $\Ext$ degree two can follow or precede
a class of $\cH^{3}(L_{\zeta},L_{\zeta'})$, so the differential no longer
vanishes on these classes. When the three shifts are consecutive, such terms
reach only the groups
between the highest and the lowest shift, which are too small to remove the
classes.

\begin{theorem}[Three consecutive shifts on $E_{0}^{8}$]
\label{thm:efourthreelevels}
Keep \Cref{lem:efourshifts}, and suppose that the shifts of the $L_{\zeta}$
lie in $\{s,s-1,s-2\}$ for an integer $s$. By their parities, for a sign
$\epsilon$, the pieces with $\prod\zeta=\epsilon$ have the shift $s-1$ and
the others the shift $s$ or $s-2$; let $T$ and $B$ be the sets of pieces of
the shifts $s$ and $s-2$. Then, for every order of the pieces, every choice
of the shifts of the $M_{i}$ and every placement of the $t_{i}$:
\begin{enumerate}[label=\textup{(\roman*)},leftmargin=2.6em,itemsep=2pt]
\item every two pieces $L_{\zeta}$, $L_{\zeta'}$ that differ in exactly three
  coordinates have shifts differing by one, and the groups
  $\cH^{3}(L_{\zeta},L_{\zeta'})$ of \Cref{thm:efourtwolevels}, taken from
  the piece of the larger shift, have dimensions adding up to $40960$;
\item $d_{E}$ maps the sum of these groups into
  $\bigoplus_{c\in T,\,e\in B}\cH^{5}(c,e)$, and on them it consists of cup
  products with a single component of $\Ext$ degree two;
\item
  \[
    \dim\Ext^{2}(E,E)\;\ge\;40960-\sum_{c\in T,\,e\in B}\dim\cH^{5}(c,e)
    \;\ge\;40960-18432=22528 .
  \]
\end{enumerate}
In particular $E$ does not meet the polarised criterion. When $T$ or $B$ is
empty this is \Cref{thm:efourtwolevels}(ii).
\end{theorem}

\begin{proof}
(i) The ratios $\zeta'_{j}/\zeta_{j}$ in the three coordinates have
product $-1$, so the two pieces have opposite parities: one has the shift
$s-1$ and the other $s$ or $s-2$. Each of the $64$ pieces of the shift $s-1$
has $28$ such partners, whose groups have dimensions adding up to $640$, as
in the proof of \Cref{thm:efourtwolevels}; this gives $64\cdot640=40960$.

(ii) Let $x\in\cH^{3}(a,b)$ with $d_{a}=d_{b}+1$. As in the proof of
\Cref{thm:efourtwolevels}(ii), a term of $d_{E}(x)$ runs along chains of
components $c\to\dots\to a$ and $b\to\dots\to e$, not both empty, and lands
in $\cH(c,e)$ in $\Ext$ degree $3+d_{c}-d_{e}$. The shifts of the pieces
$L_{\zeta}$ differ by at most two, so by \Cref{lem:efourshifts}(iii) no part
of a chain between two pieces $L_{\zeta}$ passes through a multiple. Hence a
chain $c\to\dots\to a$ that starts at a multiple $M_{i}$ is a run
$M_{i}=m_{0}\to\dots\to m_{r}\to L_{1}$ of distinct multiples, $r\le2$,
followed by pieces $L_{\zeta}$ only; let $U_{1}$ and $D_{1}$ be the numbers
of its steps up and down the value $\phi$ before $L_{1}$, so that
$d_{L_{1}}=d_{M_{i}}+U_{1}-7D_{1}$ by \Cref{lem:efourshifts}(ii) and
$U_{1}+D_{1}=r+1\le3$. Likewise a chain $b\to\dots\to e$ ending at a multiple
$M_{k}$ ends with a run from a piece $L_{2}$ with $U_{2}$ steps up and $D_{2}$
down, $U_{2}+D_{2}\le3$ and $d_{M_{k}}=d_{L_{2}}+U_{2}-7D_{2}$. Put
$L_{1}=c$ if $c$ is a piece $L_{\zeta}$ and $L_{2}=e$ if $e$ is one, and
$\delta=d_{L_{1}}-d_{L_{2}}$; then $1=d_{a}-d_{b}\le\delta\le2$.

If $c$ and $e$ are both pieces $L_{\zeta}$, each nonempty chain lowers the
shift by at least one, so $3+d_{c}-d_{e}\ge5$, while $\delta\le2$ gives at
most $5$. So the degree is $5$, $c\in T$, $e\in B$, and the two chains
together consist of one component, of $\Ext$ degree two by
\Cref{lem:efourshifts}(i): the term is a cup product of $x$ with that
component.

If $c=M_{i}$ and $e$ is a piece $L_{\zeta}$, the degree is
$3+\delta-U_{1}+7D_{1}$, and $\cH(c,e)$ sits in degree $0$ or $8$ by
\Cref{lem:efourshifts}(ii). Degree $0$ needs $U_{1}\ge3+\delta\ge4$, and
degree $8$ needs $7D_{1}-U_{1}=5-\delta\in\{3,4\}$, so $D_{1}\ge1$ and
$U_{1}\ge3$; both contradict $U_{1}+D_{1}\le3$. The case of a piece $c$ and
a multiple $e$ is the same with the second run.

If $c=M_{i}$ and $e=M_{k}$, the degree is $3+\delta-U+7D$ with
$U=U_{1}+U_{2}\le6$ and $D=D_{1}+D_{2}$. For $i=k$ it is $3$, so
$U-7D=\delta\ge1$; $D=0$ is impossible, since the values would increase from
$t_{i}$ to $p$ along the first run and from $p$ to $t_{i}$ along the second,
and $D\ge1$ gives $U\ge8$. For $i\ne k$, $\cH(c,e)$ sits in degree $0$ if
$t_{k}>t_{i}$ and $8$ if $t_{k}<t_{i}$. Degree $0$ needs $U=3+\delta+7D$,
so $D=0$ and $U\ge4$; but then the multiples of the first run have values
below $p$ and those of the second values above $p$, so they are distinct and
$U\le3$. Degree $8$ needs $7D-U=5-\delta$, so $D=1$ and $U\ge3$: the runs
have at least four steps, so some multiple lies on both. One run increases
strictly. If it is the second, its values lie in $(p,t_{k}]$, while those of
the first lie at or above $t_{i}>t_{k}$ before its step down and below $p$
after it. If it is the first, its values lie in $[t_{i},p)$, while those of
the second lie above $p$ before its step down and at or below $t_{k}<t_{i}$
after it. Either way no multiple lies on both runs, a contradiction.

(iii) By \Cref{thm:efourtwolevels}(i), which holds for every convolution, no
coboundary has a component in these groups, so the kernel of $d_{E}$ on
their sum injects into $\Ext^{2}(E,E)$, and by (ii) its dimension is at
least $40960$ minus the dimension of the targets. For pieces $c$, $e$ of the
same parity, $\cH^{5}(c,e)$ is nonzero only when they differ in two or three
coordinates, and each piece has such partners with groups $\cH^{5}$ of
dimensions adding up to $960$.\footnote{The ratios $e_{j}/c_{j}\ne1$ have
product $1$. In two coordinates they are $(-1,-1)$, with $D=16$, or
$(i,-i)$, $(-i,i)$, with $D=4$; then
$\cH^{5}=\cH^{1}\otimes\cH^{1}\otimes\cH^{3}(B_{j}\times B_{k},\cO)$ has
dimension $4D$, and the six pairs of coordinates give
$6\cdot64+12\cdot16=576$. In three coordinates they are a permutation of
$(i,i,-1)$ or of $(-i,-i,-1)$, with $D=16$, and
$\cH^{5}=(\cH^{1})^{\otimes3}\otimes H^{2}(B_{l},\cO)$ has dimension $16$;
the four triples give $4\cdot6\cdot16=384$. In four coordinates
$\cH(c,e)$ sits in degree $4$.}\enlargethispage{3pt} So the sum in (iii) is at most
$\min(|T|,|B|)\cdot960\le32\cdot960=30720$, which already gives
$\dim\Ext^{2}(E,E)\ge10240$. The dimension of $\cH^{5}(c,e)$ depends only on
the ratio $e/c$, so the sum is the weight of the cut $(T,B)$ of a weighted
Cayley graph on the group of ratios of two pieces of one parity, of order
$64$ and degree $960$. For the vector $x=\pm1$ of the cut, the weight is
$\tfrac14\bigl(64\cdot960-x^{\mathsf T}Ax\bigr)\le16\,(960-\lambda)$, with
$A$ the weighted adjacency matrix and $\lambda$ its least eigenvalue. Its
eigenvalues are the character sums $\sum_{g}\dim\cH^{5}(1,g)\,\chi(g)$. A
character is $\chi(g)=\prod_{j}g_{j}^{k_{j}}$ with $k\in(\ZZ/4)^{4}$, up to
adding $(1,1,1,1)$; put $\omega_{j}=i^{k_{j}}$, $W=\sum_{j}\omega_{j}$,
$\varrho=\sum_{j}\omega_{j}^{2}$ and $\Omega=\prod_{j}\omega_{j}$. By the
list in the footnote, the ratios that move the pair of coordinates
$\{j,l\}$ contribute $64\,\omega_{j}^{2}\omega_{l}^{2}
+32\,\Re(\omega_{j}\bbar\omega_{l})$, and those that fix only the
coordinate $o$ contribute $32\,\Re\bigl(\Omega\bbar\omega_{o}(W-\omega_{o})\bigr)$.
Summing,
\[
  \chi\bigl(\textstyle\sum_{g}\dim\cH^{5}(1,g)\,g\bigr)
  =32\varrho^{2}+16|W|^{2}-192+32\,\Re(\Omega)\bigl(|W|^{2}-4\bigr).
\]
If $\Re\Omega=0$, the sum of the $k_{j}$ is odd, so an odd number of them
are odd, $|\varrho|=2$, and the value is at least $-64$. If $\Re\Omega=1$,
an even number are odd: with none or four, $|\varrho|=4$ and the value is at
least $192$; with two, $\varrho=0$, and $W=0$ would force the $k_{j}$ to be
$0,1,2,3$, whose sum is not divisible by $4$, so $|W|^{2}\ge4$ and the value
is at least $-128$. If $\Re\Omega=-1$, the value is
$32\varrho^{2}-16|W|^{2}-64$, and $|W|^{2}\le8+2\varrho^{2}$: when all
$k_{j}$ have the same parity, $|W|\le4=|\varrho|$, and when two are odd,
$\varrho=0$ and $W$ is a sum of two of $\pm1$ and two of $\pm i$, so
$|W|^{2}\le8$. Hence $\lambda\ge-192$,\footnote{The eight values are $960$,
$384$, $192$, $96$, $-32$, $-64$, $-128$ and $-192$. The bound $18432$ is
attained by the cut $\zeta_{3}\zeta_{4}\in\{1,i\}$ against
$\zeta_{3}\zeta_{4}\in\{-1,-i\}$, so $22528$ is the best bound that this
count gives.} and the sum is at most $16\cdot1152=18432$.
\end{proof}

\Cref{fig:efourthree} draws the three shifts and the budget.

\begin{figure}[tbp]
\centering
\includegraphics[width=\textwidth]{fig_efourthree}
\caption{\Cref{thm:efourthreelevels}. Left: the pieces $L_{\zeta}$ at three
consecutive shifts, those of one parity at $s-1$ and the others split between
$T$ at $s$ and $B$ at $s-2$; a group $\cH^{3}$ between two pieces that differ
in three coordinates (teal) can meet a component of $\Ext$ degree two
(indigo) only in a group $\cH^{5}(c,e)$ with $c\in T$ and $e\in B$ (magenta).
Middle: the $40960$ classes against the targets, at most $30720$ by counting
and $18432$ by the least eigenvalue, against the $104$ that the criterion
allows. Right: the eigenvalues of the Cayley graph of the targets.}
\label{fig:efourthree}
\end{figure}

The groups of \Cref{thm:efourtwolevels} join two pieces whose shifts differ
by one. All groups between pieces whose shifts differ by three, and thirteen
kinds between pieces whose shifts differ by two, are also left alone by the
differential, and with them the argument covers every
arrangement of the shifts within five consecutive values. For pieces $a$ and
$b$ call $d_{a}-d_{b}$ the \emph{spread} of $\cH(a,b)$; its classes of degree
two have $\Ext$ degree $2+d_{a}-d_{b}$. For $\zeta,\zeta'\in\mu_{4}^{4}$ call
$\zeta'/\zeta=(\zeta'_{j}/\zeta_{j})_{j}$ their \emph{ratio}; it
\emph{moves} the coordinates where it differs from $1$, and it changes the
parity, $\prod\zeta'=-\prod\zeta$, exactly when the product of its entries is
$-1$. As before,
$D(\zeta,\zeta')=\prod_{\zeta_{j}\ne\zeta'_{j}}|\zeta_{j}-\zeta'_{j}|^{2}$.

\begin{lemma}[Groups of spread two and three]\label{lem:efourspread}
Keep \Cref{lem:efourshifts}, and let $E$ be any convolution of the
$L_{\zeta}$ and the $M_{i}$, in any order, with any shifts of the given
parities and any placement of the $t_{i}$.
\begin{enumerate}[label=\textup{(\roman*)},leftmargin=2.6em,itemsep=2pt]
\item For a class $x$ of degree two in $\cH(L_{\zeta},L_{\zeta'})$ of
  positive spread, every term of $d_{E}(x)$ has all its leaves between
  pieces $L_{\zeta''}$.
\item Let $\prod\zeta=\prod\zeta'$, let $\zeta'/\zeta$ move all four
  coordinates, one of them by $-1$, and let $d_{L_{\zeta}}=d_{L_{\zeta'}}+2$.
  The classes of degree two of $\cH(L_{\zeta},L_{\zeta'})$ form
  $\cH^{4}(L_{\zeta},L_{\zeta'})$, of dimension $D(\zeta,\zeta')\in\{64,256\}$;
  no element of degree one has a differential with a nonzero component in
  it, and $d_{E}$ vanishes on it.
\item Let $\prod\zeta=-\prod\zeta'$ and $d_{L_{\zeta}}=d_{L_{\zeta'}}+3$. The
  classes of degree two of $\cH(L_{\zeta},L_{\zeta'})$ lie in
  $\cH^{5}(L_{\zeta},L_{\zeta'})$, and $d_{E}$ vanishes on them.
\item Let $\zeta$ and $\zeta'$ differ in exactly three coordinates, with
  $d_{L_{\zeta}}=d_{L_{\zeta'}}+1$. Every term of $d_{E}$ on
  $\cH^{3}(L_{\zeta},L_{\zeta'})$ runs along chains of components
  $c\to\dots\to L_{\zeta}$ and $L_{\zeta'}\to\dots\to e$ between pieces
  $L_{\zeta''}$ that together lower the shift by one or two.
\end{enumerate}
\end{lemma}

\begin{proof}
Write $a=L_{\zeta}$, $b=L_{\zeta'}$ and $\delta=d_{a}-d_{b}$. As in the proof
of \Cref{thm:efourtwolevels}(ii), a term of $d_{E}$ on a class of degree two
in $\cH(a,b)$ runs along chains of components $c\to\dots\to a$ and
$b\to\dots\to e$, not both empty, and lands in $\cH(c,e)$ in $\Ext$ degree
$3+d_{c}-d_{e}$, which is $3$ if $c=e$.

(i) Let $L_{1}$ be the first piece $L_{\zeta''}$ of the first chain and
$L_{2}$ the last of the second, and put
$X=(d_{L_{1}}-d_{a})+(d_{b}-d_{L_{2}})$. By \Cref{lem:efourshifts}(iii),
$X\ge0$, and $X\ge4$ if a multiple lies between $L_{1}$ and $a$ or between
$b$ and $L_{2}$. Before $L_{1}$ the first chain is a run of distinct
multiples, with $U_{1}$ steps up and $D_{1}$ steps down the value $\phi$ and
$U_{1}+D_{1}\le3$; after $L_{2}$ the second chain is a run with $U_{2}$ and
$D_{2}$ steps, $U_{2}+D_{2}\le3$. By \Cref{lem:efourshifts}(ii) the term has
$\Ext$ degree $3+\delta+X-U+7D$, where $U=U_{1}+U_{2}$ and $D=D_{1}+D_{2}$.
If $c$ and $e$ are pieces, this is $3+\delta+X$; a multiple on a chain makes
it at least $8$, but $\cH^{8}(c,e)\ne0$ only for $c=e$, where the degree is
$3$. If $c=M_{i}$ and $e$ is a piece, $\cH(c,e)$ sits in degree $0$ if
$t_{i}<p$ and in degree $8$ if $t_{i}>p$. Degree $0$ needs
$U_{1}\ge3+\delta+X\ge4$. Degree $8$ needs a step down on the run from
$t_{i}$ to $p$, so $D_{1}\ge1$, and then
$U_{1}=7D_{1}-5+\delta+X\ge3$. Both contradict $U_{1}+D_{1}\le3$, and a
piece $c$ with a multiple $e$ is the same. Let $c=M_{i}$ and $e=M_{k}$. For
$i=k$ the degree is $3$, so $U-7D=\delta+X\ge1$; $D\ge1$ would give $U\ge8$,
and $D=0$ would make the values increase from $t_{i}$ to $p$ and then from
$p$ to $t_{i}$. For $i\ne k$, $\cH(c,e)$ sits in degree $0$ if $t_{k}>t_{i}$
and in degree $8$ if $t_{k}<t_{i}$. Degree $0$ needs $U=3+\delta+X+7D$, so
$D=0$ and $U\ge4$; but then the first run lies below $p$ and the second
above, so the multiples are distinct and $U\le3$. Degree $8$ needs $D\ge1$,
since with $D=0$ the values would satisfy $t_{i}<p<t_{k}$; then
$7D-U=5-\delta-X\le4$ gives $D=1$ and $U\ge3$, so the runs have at least
four steps and share a multiple, which the last paragraph of the proof of
\Cref{thm:efourthreelevels}(ii) excludes.

(ii) By the proof of \Cref{lem:efourshifts}(i), $\cH(a,b)$ is the tensor
product of four groups $H^{1}$ and sits in degree $4$, of dimension
$D(\zeta,\zeta')$; the ratio is $(-1,-1,-1,-1)$, with $D=256$, or a
permutation of $(i,-i,-1,-1)$, with $D=64$.\footnote{A ratio that moves
every coordinate has entries in $\{i,-1,-i\}$ and product $1$. With a $-1$
among them it is $(-1,-1,-1,-1)$ or has exactly two entries $-1$ and the
entries $i$ and $-i$, twelve ratios; together thirteen, and
$256+12\cdot64=1024$. The other eight ratios have all four entries in
$\{i,-i\}$ and $D=16$; (ii) makes no statement about them.}\enlargethispage{3pt} A coboundary with a component in $\cH(a,b)$ is
a sum of terms of \Cref{lem:harmonicmodel}(iii) along a path
$a\to\dots\to b$ of morphisms of degree one, one of them an element $u$ and
the others components, that lowers the shift by two. By
\Cref{lem:efourshifts}(iii) no multiple lies on it, and by
\Cref{lem:efourshifts}(i) each step between distinct pieces lowers the shift
by at least one. So the path consists of a component of $\Ext$ degree three
from $a$ to $b$, which lies in $\cH^{3}(a,b)=0$, and a class of
$\cH^{1}(a,a)$ or $\cH^{1}(b,b)$; or of two steps $a\to c\to b$ of $\Ext$
degree two, possibly with a class of $H^{1}(A,\cO_{A})$ at one of the three
pieces. The output has degree one on every surface. Without the class of
$H^{1}$ the two steps have degrees adding up to one on every surface, so they
move disjoint pairs of coordinates and have degree zero elsewhere; a step of
$\Ext$ degree two changes the parity, so its two moved entries have product
$-1$ and are $(i,i)$ or $(-i,-i)$, and $\zeta'/\zeta$ has no entry $-1$.
With the class of $H^{1}$, \Cref{lem:degreedrop} gives a drop of one on a
surface where all three leaves have positive degree, and the output has
degree at least two there.

Now let a term of $d_{E}$ on $\cH^{4}(a,b)$ carry the components
$y_{1},\dots,y_{m}$, $m\ge1$, between pieces by (i), and land in
$\cH(c,e)$; let $K$ be the set of coordinates fixed by the ratio of $c$ and
$e$, and $X\ge m$ the drop of the chains. The output has $\Ext$ degree
$5+X$. Let $n_{j}$ be the number of the $y_{i}$ of positive degree on
$B_{j}$; the class $x$ has degree one there, so \Cref{lem:degreedrop} bounds
the drop on $B_{j}$ by $\max(0,n_{j}-1)$. For $j\notin K$ the output has
degree $1$ on $B_{j}$, so the degrees of the $y_{i}$ there add up to at most
$\max(0,n_{j}-1)$, which forces $n_{j}=0$. For $j\in K$ the $y_{i}$ undo the
entry $\zeta'_{j}/\zeta_{j}\ne1$, so $n_{j}\ge1$, and the output has degree
at least $1+n_{j}-(n_{j}-1)=2$ there, hence exactly $2$. So
$5+X=(4-|K|)+2|K|$, that is $X=|K|-1$. The components change the parity
$X$ times modulo two, and together their ratio is $\zeta/\zeta'$ on $K$ and
$1$ elsewhere, so
\[
  \prod_{j\in K}\zeta'_{j}/\zeta_{j}=(-1)^{|K|-1}.
\]
For the ratio $(-1,-1,-1,-1)$ the left side is $(-1)^{|K|}$. For a
permutation of $(i,-i,-1,-1)$ the left side is real only if $K$ contains
both or neither of the coordinates moved by $\pm i$, and then it is again
$(-1)^{|K|}$. So no term is nonzero.

(iii) The classes of degree two have $\Ext$ degree $5$. By (i) a term of $d_{E}$ on them lands
in $\Ext$ degree $6+X$, with $X\ge1$ the drop of the chains between pieces.
For $X\ge2$ this is at least $8$, which needs $c=e$ and degree $3$. For $X=1$
the pieces $c$, $e$ have shifts differing by four, hence the same parity, so
their ratio moves at least two coordinates and $\cH(c,e)$ has no degree above
$6$.

(iv) By (i) a term lands in $\Ext$ degree $4+X$, with $X\ge1$ the drop of the
chains between pieces. For $X\ge4$ the degree is at least $8$, which needs
$c=e$. For $X=3$ the degree is $7$, which needs a ratio moving one
coordinate, hence opposite parities, while $d_{c}-d_{e}=4$ is even.
\end{proof}

The groups of (ii) are reached by no coboundary and left by no term in every
convolution, so they survive in $\Ext^{2}(E,E)$ as soon as the shifts give
them the spread two. That happens along every edge of a Cayley graph that
crosses between two shift values of one parity, and a cut of that graph is
never light.

\begin{lemma}[Cuts of a Cayley graph]\label{lem:cayleycut}
Let $G$ be a finite abelian group, $S\subset G\setminus\{0\}$ a generating
set and $w\colon S\to\RR_{>0}$ a weight, extended by $0$ to $G$. For every
subset $T\subset G$ with $T\ne\emptyset,G$,
\[
  \sum_{c\in T,\ e\notin T}w(e-c)\;\ge\;\min_{H}\,|H|\cdot w(S\setminus H),
\]
the minimum over the proper subgroups $H$ of $G$. For the group $G_{0}$ of
the $64$ ratios of pieces of one parity, with $S$ the thirteen ratios of
\Cref{lem:efourspread}(ii) and $w=D$, the minimum is $1024$, attained only at
$H=0$, so every such cut weighs at least $1024$.
\end{lemma}

\begin{proof}
The left side is $\sum_{s\in S}w(s)\,|(T+s)\setminus T|$. Let $H$ be the
stabiliser of $T$, a proper subgroup since $T\ne\emptyset,G$. Then $T$ and
$T+s$ are unions of cosets of $H$, and for $s\notin H$ the set
$(T+s)\setminus T$ is nonempty, since $|T+s|=|T|$, so it has at least $|H|$
elements. This gives the inequality. In $G_{0}$, written additively as the
quadruples in $(\ZZ/4)^{4}$ of sum $0$, the ratios of type $(i,-i,-1,-1)$
generate, since the differences of two of them give every $e_{j}-e_{k}$. The
weights are $256$ once and $64$ otherwise, and $|S\cap H|\le|H|-1$. So
$|H|\,w(S\setminus H)$ is $1024$ for $H=0$, at least
$|H|\,(1024-256-64(|H|-2))$, which is $1536$, $2560$ and $3072$, for
$|H|=2,4,8$, and at least $64|H|\ge1024$ for $|H|\ge16$, because some
generator lies outside $H$.\footnote{The orders of the subgroups of
$G_{0}\cong(\ZZ/4)^{3}$ are powers of two. A cut at one vertex weighs exactly
$1024$, so the bound is sharp.}
\end{proof}

\begin{theorem}[One parity at two or three values]\label{thm:efourspreadtwo}
Keep \Cref{lem:efourshifts}, and suppose that the shifts of the pieces
$L_{\zeta}$ with $\prod\zeta=\epsilon$ take exactly two values $s$ and $s-2$,
or exactly three values $s+2$, $s$ and $s-2$, for a sign $\epsilon$ and an
integer $s$. Then, for every order of the pieces, all shifts of the other
pieces, every choice of the shifts of the $M_{i}$ and every placement of the
$t_{i}$,
\[
  \dim\Ext^{2}(E,E)\;\ge\;1024 .
\]
In particular $E$ does not meet the polarised criterion.
\end{theorem}

\begin{proof}
Let $T$ be the set of pieces of parity $\epsilon$ at the shift $s$; it is
neither empty nor all of them. Two pieces of parity $\epsilon$, one in $T$
and the other not, have shifts differing by exactly two. When their ratio is
one of the thirteen of \Cref{lem:efourspread}(ii), the group $\cH^{4}$ from
the piece of the larger shift to the other is reached by no coboundary and
killed by no term, so its classes inject into $\Ext^{2}(E,E)$, and those of
different pairs lie in different summands $\cH(a,b)$. Their dimensions add up
to the weight of the cut $(T,\complement T)$ of the Cayley graph of
\Cref{lem:cayleycut}, since the pieces of parity $\epsilon$ form a coset of
$G_{0}$; so they add up to at least $1024$.
\end{proof}

\begin{theorem}[Shifts three apart]\label{thm:efourspreadthree}
Keep \Cref{lem:efourshifts}, and suppose that the pieces $L_{\zeta}$ with
$\prod\zeta=\epsilon$ all have the shift $s$ and the others have shifts in
$\{s-3,s+3\}$. Then, for every order, every choice of the shifts of the
$M_{i}$ and every placement of the $t_{i}$,
\[
  \dim\Ext^{2}(E,E)\;\ge\;64\cdot1072-128\cdot8=67584 .
\]
\end{theorem}

\begin{proof}
Every piece $W$ of parity $-\epsilon$ and every piece $V$ of parity $\epsilon$
have shifts differing by three. By \Cref{lem:efourspread}(iii) the classes
of degree two of the group from the higher to the lower of the two lie in
$\cH^{5}$, and $d_{E}$ vanishes on them. For each $W$ these groups have
dimensions adding up to $1072$.\footnote{Let the ratio move $k$
coordinates. For $k=1$ it is $-1$ in one coordinate, $D=4$, and $\cH^{5}$
is $H^{1}$ tensored with the part of degree four of $H^{*}(\cO)$ on the
other three surfaces, of dimension $\binom{6}{4}=15$: four partners with
$60$ each. For $k=2$ it is $(i,i)$ or $(-i,-i)$, $D=4$, and the part of
degree three on the other two surfaces has dimension $\binom43=4$: twelve
partners with $16$ each. For $k=3$ the part of degree two on the last surface
has dimension one, and the $28$ partners give $640$ as in
\Cref{thm:efourtwolevels}. For $k=4$ the group sits in degree $4$. So
$240+192+640=1072$.} A coboundary with a component in such a group runs
along a path of morphisms of degree one that lowers the shift by three;
no multiple lies on it, by \Cref{lem:efourshifts}(iii), and a step between
distinct pieces lowers the shift by at least one, while no piece has a shift
strictly between those of $W$ and $V$. So the path is the component between
$W$ and $V$ together with a class of $H^{1}(A,\cO_{A})$ at one of them, and
the components of all coboundaries in these groups come from
$\bigoplus_{a}\cH^{1}(a,a)$ over the $128$ pieces, of dimension
$128\cdot8=1024$. The $64\cdot1072=68608$ classes are cocycles, and at most
$1024$ dimensions of them are coboundaries.
\end{proof}

When one parity sits at a single shift, the groups $\cH^{3}$ of
\Cref{thm:efourtwolevels} survive next to it unless pieces of the other
parity crowd it from both sides.

\begin{proposition}[A gap of four]\label{prop:efourgap}
Keep \Cref{lem:efourshifts}, let the pieces $L_{\zeta}$ with
$\prod\zeta=\epsilon$ all have the shift $y$, and let $T$ be the set of
pieces of parity $-\epsilon$ at the shift $y+1$. If $T$ is not empty and no
piece has the shift $y-1$ or $y+3$, then for every order, every choice of
the shifts of the $M_{i}$ and every placement of the $t_{i}$, the groups
$\cH^{3}(W,V)$ with $W\in T$ and $V$ differing from $W$ in exactly three
coordinates are reached by no coboundary and killed by no term, and
\[
  \dim\Ext^{2}(E,E)\;\ge\;640\,|T|\;\ge\;640 .
\]
The same holds with $y+1$, $y-1$, $y+3$ replaced by $y-1$, $y+1$, $y-3$ and
the groups $\cH^{3}(V,W)$.
\end{proposition}

\begin{proof}
By \Cref{thm:efourtwolevels}(i), which holds for every convolution, no
coboundary reaches these groups. By \Cref{lem:efourspread}(iv) a term of
$d_{E}$ on $\cH^{3}(W,V)$ runs along chains of components between pieces
into $W$ and out of $V$ that lower the shift by one or two in all; the
first would pass through a piece at the shift $y+2$ or $y+3$, the second
through one at $y-1$ or $y-2$. The shifts $y\pm2$ have parity $\epsilon$,
whose pieces are all at $y$, and $y-1$, $y+3$ are excluded. Each $W\in T$
has $28$ partners $V$, all at the shift $y$, with groups of dimensions adding
up to $640$. The second statement has the same proof, with the chains into
$V$ and out of $W$ in place of those into $W$ and out of $V$.
\end{proof}

\begin{corollary}[Five consecutive shifts]\label{cor:efourfive}
Keep \Cref{lem:efourshifts}. If the shifts of the $128$ pieces $L_{\zeta}$
lie in five consecutive integers, then, for every order, every choice of the
shifts of the $M_{i}$ and every placement of the $t_{i}$,
$\dim\Ext^{2}(E,E)\ge640$; no such convolution meets the polarised
criterion.
\end{corollary}

\begin{proof}
Let the values be $s,\dots,s-4$, with one parity at values in
$\{s,s-2,s-4\}$ and the other at values in $\{s-1,s-3\}$. If the second
parity takes both its values, or the first takes two values two apart or all
three, \Cref{thm:efourspreadtwo} applies. Otherwise the second parity sits at
one value $y$ and the first at one value or at $\{s,s-4\}$. One value each,
at an odd distance at most three: \Cref{thm:efourtwolevels} or
\Cref{thm:efourspreadthree}. The first at $\{s,s-4\}$: for $y=s-1$ the
pieces at $y+1=s$ are not empty and $y-1$, $y+3$ are empty, and for $y=s-3$
the pieces at $y-1=s-4$ are not empty and $y+1$, $y-3$ are empty, so
\Cref{prop:efourgap} applies. The least of the bounds is $640$.
\end{proof}

Within six consecutive values two arrangements are left, up to exchanging
the parities: each parity at one shift, five apart, and the parities at
$\{s,s-4\}$ and $\{s-1,s-5\}$. The first is settled by the diagonal, whose
classes at a shift with no other piece nearby can be reached only through
all three multiples, and the second by the groups of
\Cref{thm:efourtwolevels} once more.

\begin{lemma}[The diagonal at a lonely shift]\label{lem:efourlonely}
Keep \Cref{lem:efourshifts}, and let the pieces $L_{\zeta}$ with
$\prod\zeta=\epsilon$ all have the shift $s$ and the others the shift $s-g$,
for a sign $\epsilon$ and an odd integer $g\ge5$. Let $y\in\cH^{2}(a,a)$ be
a diagonal class at a piece $a$, and let a term of $d_{E}(y)$ in
\Cref{lem:harmonicmodel}(iii) be nonzero. Then $a$ is the only piece on its
path, the path passes through $M_{1}$, $M_{2}$ and $M_{3}$, and the shifts
of $a$, $M_{1}$, $M_{2}$, $M_{3}$ are pairwise incongruent modulo $4$.
\end{lemma}

\begin{proof}
Write the path as $c\to\dots\to a\to\dots\to e$, with the chain of
components into $a$ and the chain out of $a$ not both empty, so that
$c\ne e$, since the components of a convolution go up the order of the
pieces. Let $L_{1}$ and $L_{2}$ be the first and the last piece on it and
$X=d_{L_{1}}-d_{L_{2}}$; before $L_{1}$ and after $L_{2}$ the path is a run
of distinct multiples, with $U$ steps up and $D$ steps down the value $\phi$
in all. As in the proof of \Cref{lem:efourspread}(i), now with $\delta=0$,
the term has $\Ext$ degree $3+X-U+7D$ by \Cref{lem:efourshifts}(ii). A
sequence of components from one piece to another lowers the shift by
\Cref{lem:efourshifts}(iii), and the pieces have the shifts $s$ and $s-g$;
so $X\in\{0,g\}$, and $X=0$ exactly when $a$ is the only piece on the path.

If $c$ and $e$ are pieces, then $U=D=0$ and $X=g$, and $\cH^{3+g}(c,e)=0$
since $3+g\ge8$ and $c\ne e$. If $c$ is a multiple and $e$ a piece, then
$\cH(c,e)$ sits in degree $0$ or $8$. Degree $0$ needs $U=3+X+7D$, so
$X=D=0$ and $U=3$. Degree $8$ needs $7D-U=5-X$: for $X=0$ this gives $D=1$
and $U=2$; for $X=g$ it gives $D=0$ and $U=g-5$, which is at least one
because the run before $L_{1}$ is not empty and at most three, so $g=7$; but
then the path lowers the shift by seven from $L_{1}$ to $e$, by a single
component, of $\Ext$ degree eight between distinct pieces, or by a run
through multiples, and neither is possible.\footnote{Between distinct pieces
$\cH^{8}=0$. A run from a piece through $m\le3$ distinct multiples to a piece
has $m+1$ steps, at least one up and one down, and lowers the shift by
$7D'-U'$ with $U'+D'=m+1$: by $6$, $5$, $4$ for $D'=1$, by $13$, $12$ for
$D'=2$ and by $20$ for $D'=3$, never by seven.} A piece $c$ with a multiple $e$ is the same
with the runs exchanged. If $c$ and $e$ are multiples, both runs are
nonempty and together pass through at most three multiples, so $U+D\le3$
and $U+D\ge2$. Degree $0$ needs $U=3+X+7D$, so $X=D=0$ and $U=3$. Degree $8$
needs $7D-U=5-X$; $D=0$ gives $U=X-5\ge2$, so $X=g=7$, which is excluded as
before; $D=1$ gives $U=2+X\le2$, so $X=0$ and $U=2$; $D\ge2$ gives $U\ge9$.

So in every case $X=0$ and $U+D=3$: the path consists of $a$ and the three
multiples, with three steps. A step up raises the shift by one and a step
down lowers it by seven, so every step adds one modulo $4$, and the four
nodes of the path have the shifts $d_{c},d_{c}+1,d_{c}+2,d_{c}+3$ modulo $4$
in the order of the path.\footnote{With the values $t_{1}<t_{2}<t_{3}$ and
$p$, the path with $U=3$ visits the four nodes in the order of their values;
a path with $U=2$ and $D=1$ ending in degree eight visits the largest $j$
values in increasing order and then the others, $j=1,2,3$. So each placement
has exactly four such paths, and for each the shifts of the multiples are
determined by $d_{a}$.}
\end{proof}

\begin{proposition}[Single shifts]\label{prop:efoursingle}
Keep \Cref{lem:efourshifts}, and suppose that the pieces $L_{\zeta}$ with
$\prod\zeta=\epsilon$ all have the shift $s$ and the others all have the
shift $s-g$, for a sign $\epsilon$ and an odd integer $g$. Then, for every
order, every choice of the shifts of the $M_{i}$ and every placement of the
$t_{i}$,
\[
  \dim\Ext^{2}(E,E)\;\ge\;64\cdot28=1792 .
\]
\end{proposition}

\begin{proof}
Exchanging the parities, we may take $g>0$. For $g=1$ and $g=3$ this is
\Cref{thm:efourtwolevels} and \Cref{thm:efourspreadthree}. Let $g\ge5$. If a
term of $d_{E}$ were nonzero on a diagonal class at a piece of shift $s$ and
another on a diagonal class at a piece of shift $s-g$, then by
\Cref{lem:efourlonely} the shifts of the three multiples would be pairwise
incongruent modulo $4$ and miss both the residue of $s$ and that of $s-g$,
which differ, since $g$ is odd. So at one of the two shifts $d_{E}$ vanishes
on the $28$ diagonal classes of each of the $64$ pieces there, and by
\Cref{prop:nothingenters}, which holds for $n=4$, these $1792$ classes
inject into $\Ext^{2}(E,E)$.\footnote{The proof does not look at the classes
between the two parities. For $g=5$ these are $64\cdot16=1024$ classes of
degree two in groups $\cH^{7}$, on which $d_{E}$ vanishes by
\Cref{lem:efourspread}(i); but the coboundaries of the loops in
$\cH^{1}(a,a)$, and, when two multiples have suitable shifts, the products
along runs $W\to M_{i}\to M_{j}\to V$, reach them, and counting alone gives
nothing there.}
\end{proof}

\begin{proposition}[Interleaved shifts]\label{prop:efourpairs}
Keep \Cref{lem:efourshifts}, and suppose that the shifts of the pieces
$L_{\zeta}$ with $\prod\zeta=\epsilon$ take exactly the two values $s$ and
$s-4$, and those of the others exactly the two values $s-1$ and $s-5$. Then,
for every order, every choice of the shifts of the $M_{i}$ and every
placement of the $t_{i}$, $\dim\Ext^{2}(E,E)\ge640$.
\end{proposition}

\begin{proof}
Let $A$ and $A'$ be the pieces of parity $\epsilon$ at the shifts $s$ and
$s-4$, and $B$ and $B'$ those of parity $-\epsilon$ at $s-1$ and $s-5$; all
four are nonempty. Let $W$ and $V$ differ in exactly three coordinates, with
$d_{W}=d_{V}+1$, that is $W\in A$ and $V\in B$, or $W\in A'$ and $V\in B'$.
By \Cref{thm:efourtwolevels}(i) no coboundary reaches $\cH^{3}(W,V)$, and by
\Cref{lem:efourspread}(iv) a term of $d_{E}$ on it passes through a piece at
the shift $d_{W}+1$ or $d_{W}+2$ or at $d_{V}-1$ or $d_{V}-2$. These are
$s+1$, $s+2$, $s-2$, $s-3$ for $W\in A$ and $s-3$, $s-2$, $s-6$, $s-7$ for
$W\in A'$, and no piece has these shifts. So the groups inject into
$\Ext^{2}(E,E)$.

Let $\Gamma$ be the group of the ratios of two pieces, of order $128$,
which acts simply transitively on the pieces, and $S_{3}\subset\Gamma$ the
$28$ ratios that move exactly three coordinates and change the parity, with
the weight $D$, which adds up to $640$ (\Cref{thm:efourtwolevels}). With
$X=A\cup B'$ the pairs above are exactly the edges of the Cayley graph of
$(\Gamma,S_{3})$ between $X$ and its complement $A'\cup B$; so the
dimensions add up to the weight of that cut, and $X$ is neither empty nor
everything. Written additively, as the quadruples in $(\ZZ/4)^{4}$ of even
sum, the differences of two ratios of type $(i,-i,-1,1)$ give every
$e_{j}-e_{k}$, which generate the quadruples of sum $0$ modulo $4$, and the
elements of $S_{3}$ have sum $2$ modulo $4$; so $S_{3}$ generates $\Gamma$.
In \Cref{lem:cayleycut},
$H=0$ gives $640$; for $|H|=2,4,8,16$ at most $|H|-1$ elements of $S_{3}$ lie
in $H$, at most four of them of weight $64$, so $|H|\,w(S_{3}\setminus H)$ is
at least $2\cdot576$, $4\cdot448$, $8\cdot336$ and $16\cdot208$; for
$|H|=32$ some $r\in S_{3}$ lies outside $H$, and so does $-r$, which gives
weight at least $32$, or $64$ if $r=-r$; and for $|H|=64$ the weight outside
$H$ is at least $16$. So every such cut weighs at least $640$.\footnote{The
subgroups of $\Gamma\cong(\ZZ/4)^{3}\times\ZZ/2$ have orders that are powers
of two. A cut at a single vertex weighs exactly $640$, so the bound is
sharp.}
\end{proof}

\begin{corollary}[Six consecutive shifts]\label{cor:efoursix}
Keep \Cref{lem:efourshifts}. If the shifts of the $128$ pieces $L_{\zeta}$
lie in six consecutive integers, then, for every order, every choice of the
shifts of the $M_{i}$ and every placement of the $t_{i}$,
$\dim\Ext^{2}(E,E)\ge640$; no such convolution meets the polarised
criterion.
\end{corollary}

\begin{proof}
By \Cref{cor:efourfive} we may assume that the shifts take both values $s$
and $s-5$ of the window $s,\dots,s-5$; let the parity at $s$ have its values
in $\{s,s-2,s-4\}$ and the other in $\{s-1,s-3,s-5\}$. If one of them takes
two values two apart, or three, \Cref{thm:efourspreadtwo} applies.
Otherwise the first takes $\{s\}$ or $\{s,s-4\}$ and the second $\{s-5\}$ or
$\{s-1,s-5\}$. For $\{s\}$ and $\{s-5\}$ this is \Cref{prop:efoursingle},
and for $\{s,s-4\}$ and $\{s-1,s-5\}$ it is \Cref{prop:efourpairs}. For
$\{s\}$ and $\{s-1,s-5\}$, \Cref{prop:efourgap} applies in its second form
with $y=s$: the shift $s-1$ is taken and $s+1$, $s-3$ are not. For
$\{s,s-4\}$ and $\{s-5\}$ it applies in its first form with $y=s-5$: the
shift $s-4$ is taken and $s-6$, $s-2$ are not. The least of the bounds is
$640$.
\end{proof}

For the arrangements outside these results the diagonal still carries a
large part of the obstruction.

\begin{proposition}[Cup products on the diagonal at $n=4$]
\label{prop:efourcupkernel}
Keep \Cref{lem:efourshifts}, so that $D^{2}$ has dimension $131\cdot28=3668$.
For every convolution the cup products vanish on the $84$ diagonal classes of
the $M_{i}$, and on the $3584$ diagonal classes of the $L_{\zeta}$ they have
a kernel of dimension at least $400$. So their kernel on $D^{2}$ has
dimension at least $484$, and a convolution meeting the polarised criterion
needs products of length at least three that remove at least $380$ diagonal
classes.
\end{proposition}

\begin{proof}
As in the proof of \Cref{prop:cupkernel}, the cup products in $\cH(a,b)$
are $(y_{a}-y_{b})\wedge\delta^{H}_{ab}$, they vanish on the classes of the
$M_{i}$, and for $L_{\zeta}$, $L_{\zeta'}$ and a type $\tau$ of degree two on
the four surfaces they vanish on the part of type $\tau$ unless $\zeta'/\zeta$
fixes the coordinates in the support of $\tau$. So a class whose part of type
$\tau$ depends only on the coordinates of $\zeta$ in the support of $\tau$
has vanishing cup products. The four types $(2,0,0,0)$ have dimension one and
four values of the coordinate, the six types $(1,1,0,0)$ have dimension four
and sixteen values of the two coordinates: $4\cdot4+6\cdot16\cdot4=400$. By
\Cref{prop:nothingenters}, which holds for $n=4$, the kernel of $d_{E}$ on
$D^{2}$ injects into $\Ext^{2}(E,E)$, and it is at most $104$ only if the
longer products remove $380$ dimensions of these $484$.
\end{proof}

\Cref{fig:efourspread} and \Cref{fig:efoursix} draw the
arrangements and the bounds.

\begin{figure}[tbp]
\centering
\includegraphics[width=\textwidth]{fig_efourspread}
\caption{The shift arrangements of the pieces $L_{\zeta}$ on $E_{0}^{8}$.
Left: one parity at two values two apart; a piece $c$ is joined to the
thirteen pieces whose ratio moves every coordinate, one by $-1$, and the
groups $\cH^{4}$ of the edges crossing between the values survive, weighing
at least $1024$ (\Cref{lem:efourspread}(ii), \Cref{lem:cayleycut}). Middle:
the four kinds of arrangement within five consecutive values, with the
result that excludes each and its bound on $\dim\Ext^{2}$. Right: the
diagonal at $n=4$; the cup products leave at least $484$ of the $3668$
classes, and the criterion allows $104$ (\Cref{prop:efourcupkernel}).}
\label{fig:efourspread}
\end{figure}

\begin{figure}[tbp]
\centering
\includegraphics[width=\textwidth]{fig_efoursix}
\caption{The arrangements within six consecutive values on $E_{0}^{8}$.
Left: the parities at single shifts $s$ and $s-g$, $g$ odd and at least five;
a term on a diagonal class at a piece $a$ runs through $M_{1}$, $M_{2}$,
$M_{3}$ and no other piece, and since each step adds one to the shift modulo
$4$ the four nodes have distinct residues, so the multiples serve at most one
of the two shifts and the $1792$ diagonal classes at the other survive
(\Cref{lem:efourlonely}, \Cref{prop:efoursingle}). Middle: the interleaved
shifts; the groups $\cH^{3}$ between $A$ and $B$ and between $A'$ and $B'$
survive and form a cut of the Cayley graph of the $28$ ratios that move
three coordinates, of weight at least $640$ (\Cref{prop:efourpairs}).
Right: the four arrangements left within six values by the results for five
values, each with its bound on $\dim\Ext^{2}$ (\Cref{cor:efoursix}).}
\label{fig:efoursix}
\end{figure}

\begin{remark}[What the route gives]\label{rem:convolutionsopen}
At $n=3$ the route through convolutions of these pieces is decided. The
diagonal obstruction of \Cref{thm:esixdiagonal} leaves room only for the
fourfold products along one path, and the explicit convolution of
\Cref{prop:explicitconvolution} solves the Maurer--Cartan equation. The
obstruction that decides the route lies off the diagonal, in classes that no
product of any length can reach: those between pieces that differ in every coordinate and have
adjacent shifts (\Cref{lem:allthree}), and, at pieces whose shifts are not
adjacent to the shift of the paths, the classes of type $(1,1,0)$. No order
and no choice of shifts or twists avoids both, and
\Cref{thm:noconvolution} excludes every convolution of the $35$ pieces; in
the explicit convolution the first kind alone gives $2560$ classes against
the $57$ allowed. In each placement exactly two of the eight sign patterns
carry products, and the two are exchanged by a symmetry of the
construction.\footnote{The automorphism of $A$ that is $-1$ on the coordinate
$z_{n+1}$ and the identity on the others changes $\alpha$ into $-\alpha$,
fixes $\theta$ and permutes the $L_{\zeta}$, changing $\zeta_{1}$ into
$-\zeta_{1}$; composed with the shift $[1]$ it changes every $\sigma_{i}$
into $-\sigma_{i}$ and keeps everything else.} The member $E_{0}^{6}$ with
this Weil structure has trivial discriminant, where the Weil classes are
algebraic on every member by Markman's theorem \cite{Mar25a,Mar25b}, so a
convolution meeting the criterion here would have given a second proof of a
known case rather than a new one. At $n=4$ the classes between pieces that
differ in three coordinates and have adjacent shifts are again reached by
nothing, in every convolution, and when the pieces $L_{\zeta}$ occupy two
shifts, as in the explicit convolution, no term of $d_{E}$ leaves them, so
that $\dim\Ext^{2}(E,E)\ge40960$ against the $104$ allowed
(\Cref{thm:efourtwolevels}). With three consecutive shifts the
differential reaches these classes, but only through cup products into the
groups between the highest and the lowest shift, and at least $22528$ of
them survive (\Cref{thm:efourthreelevels}). The groups of spread three, and
those of spread two whose ratio moves every coordinate, one of them by $-1$,
are untouched as well, in every convolution
(\Cref{lem:efourspread}): the latter cross every cut between
two shift values of one parity, and weigh at least $1024$ there
(\Cref{thm:efourspreadtwo}); with one parity at a single shift, those of
spread three, or of spread one beside a gap, survive in their turn
(\Cref{thm:efourspreadthree}, \Cref{prop:efourgap}). Together they exclude
every convolution whose pieces $L_{\zeta}$ have their shifts within five
consecutive values (\Cref{cor:efourfive}). The two arrangements left within
six values fall as well: with each parity at a single shift, at any odd
distance, the diagonal classes at one of the two shifts can be reached only
through all three multiples, and the multiples serve at most one of them
(\Cref{lem:efourlonely}, \Cref{prop:efoursingle}); the interleaved shifts
leave the groups of \Cref{thm:efourtwolevels} on a cut
(\Cref{prop:efourpairs}). So every convolution whose shifts lie in six
consecutive values fails (\Cref{cor:efoursix}). For wider arrangements the
cup products still leave $484$ diagonal classes, so products of length at
least three would have to remove $380$ (\Cref{prop:efourcupkernel}), and
products of length three and four among the $L_{\zeta}$ pass every test of
\Cref{thm:esixdiagonal}. Whether some convolution of the $131$ pieces whose
shifts do not lie in six consecutive values, and whose parities do not each
sit at one shift, meets the criterion at $n=4$ is open. Markman's theorem
does not reach the eightfolds of Weil type, so such a convolution would give
a new case.
\Cref{fig:convolutions} shows the excess count and the
surviving path, and \Cref{fig:diagonalbudget} the budget of diagonal
classes.
\end{remark}
